% !TeX program = lualatex
% !TeX encoding = utf8
% !TeX spellcheck = uk_UA
% !TeX root = ../main.tex

%=========================================================
\chapter{Елементи фізики плазми}\label{\currfilebase}\hypertarget{PlasmaPhysics}{}
\graphicspath{{\currfilebase/Pictures}}
%=========================================================

У~цьому розділі використовується гаусівська (СГС) система одиниць. Тому в~законі Лоренца входить множник $1/c$, рівняння Пуассона має вигляд $\nabla^2\varphi=-4\pi\rho$, а~густина енергії магнітного поля дорівнює $B^2/(8\pi)$. Температура $T$ всюди означає термодинамічну температуру, тому теплова енергія записується як $k_{\mathrm B}T$. У~плазмофізиці температуру часто задають безпосередньо в~еВ; тоді $k_{\mathrm B}T$ замінюється на $T$ в~енергетичних формулах. Відповідність між температурою та енергією має вигляд
% ~~~~~~~~~~~~~~~~~~~~~~~~~~~~~~~~~~~~~~~~~~~~~~~~~~~~~~~~
\begin{equation}
	\frac{1~\mathrm{eV}}{k_{\mathrm B}}
	\simeq 1.160\times10^4~\mathrm{K},
\end{equation}
% ~~~~~~~~~~~~~~~~~~~~~~~~~~~~~~~~~~~~~~~~~~~~~~~~~~~~~~~~
тобто $1~\mathrm{eV}$ відповідає приблизно $1.16\times10^4$~К.

%% --------------------------------------------------------
\section{Плазма. Основні поняття}\label{sec:plasma-basics}\hypertarget{plasma-basics}{}
%% --------------------------------------------------------

%% --------------------------------------------------------
\subsection{Що таке плазма}
%% --------------------------------------------------------

\emphx{Плазма} (від грец. \emphx{plasma}~--- виліплений, оформлений)~--- це частково або повністю іонізований газ, в~якому густини позитивних і~негативних зарядів практично однакові. Термін \enquote{плазма} запропонований у~1929~р. І.~Ленгмюром.

Плазмою називають іонізований газ, у~якому заряджені частинки можуть вільно переміщуватися на макроскопічних масштабах, а~кулонівська взаємодія приводить до колективної поведінки великої кількості частинок. Плазма може бути як майже повністю, так і~лише частково іонізованою. Для характеристики ступеня іонізації введемо
% ~~~~~~~~~~~~~~~~~~~~~~~~~~~~~~~~~~~~~~~~~~~~~~~~~~~~~~~~
\begin{equation}
	\alpha = \frac{n_i}{n_i+n_n},
\end{equation}
% ~~~~~~~~~~~~~~~~~~~~~~~~~~~~~~~~~~~~~~~~~~~~~~~~~~~~~~~~
де $n_i$ і~$n_n$~--- концентрації іонів та нейтральних частинок. У~плазмі можуть одночасно бути присутніми кілька сортів іонів, електрони та нейтральні атоми або молекули.

Серед найважливіших властивостей плазми потрібно вказати те, що її механічні властивості значною мірою диктуються електродинамікою. Плазма не тільки рухається в~зовнішньому електромагнітному полі, а~й~сама створює макроскопічні електричні та магнітні поля.

Головна відмінність плазми від звичайного нейтрального газу полягає не просто в~наявності заряджених частинок. Вона визначається далекодійною кулонівською взаємодією. У~нейтральному газі основний внесок у~зміну імпульсу частинки дають короткочасні зіткнення. У~плазмі кожна заряджена частинка одночасно взаємодіє з~великою кількістю інших частинок. Тому важливими стають \emph{колективні ефекти}: екранування зарядів, плазмові коливання, хвилі, струми та взаємодія плазми з~електромагнітним полем.

Плазма виникає різними способами: термічною іонізацією, електричним розрядом, фотоіонізацією, бомбардуванням частинками, ударною іонізацією. До природної плазми належать зоряна речовина, сонячна корона, іоносфера та магнітосфера Землі. Лабораторна плазма використовується, зокрема, в~газових розрядах, плазмових технологіях, джерелах світла та дослідженнях керованого термоядерного синтезу.

Плазма називається \emphx{низькотемпературною}, якщо її температура $T < 10^5$~К, і~\emphx{високотемпературною}, якщо $T > 10^6$--$10^8$~К. Високотемпературна плазма використовується для здійснення реакцій ядерного синтезу (\emphx{керованого термоядерного синтезу}).

%% --------------------------------------------------------
\subsection{Двокомпонентність та температура плазми}
%% --------------------------------------------------------

У~плазмі електрони та іони можуть мати різні середні кінетичні енергії. Тому, взагалі кажучи, треба вводити дві температури:
% ~~~~~~~~~~~~~~~~~~~~~~~~~~~~~~~~~~~~~~~~~~~~~~~~~~~~~~~~
\begin{equation}
	T_e,\qquad T_i.
\end{equation}
% ~~~~~~~~~~~~~~~~~~~~~~~~~~~~~~~~~~~~~~~~~~~~~~~~~~~~~~~~
Тут і~далі $T$ є термодинамічною температурою. В~енергетичних формулах входить комбінація $k_{\mathrm B}T$.

У~газорозрядній плазмі енергія від зовнішнього джерела насамперед передається електронам. Через велику різницю мас
\begin{equation*}
	m_i\gg m_e
\end{equation*}
електрон при зіткненні з~важким іоном передає йому лише малу частину своєї кінетичної енергії. Тому часто
% ~~~~~~~~~~~~~~~~~~~~~~~~~~~~~~~~~~~~~~~~~~~~~~~~~~~~~~~~
\begin{equation}
	T_e>T_i.
\end{equation}
% ~~~~~~~~~~~~~~~~~~~~~~~~~~~~~~~~~~~~~~~~~~~~~~~~~~~~~~~~
За достатньо великої густини та тривалого часу взаємодії електрон-іонні зіткнення приводять до теплообміну між компонентами, і~в~стані повної термодинамічної рівноваги
\begin{equation*}
	T_e=T_i=T.
\end{equation*}

Отже, слово \enquote{температура плазми} без додаткового уточнення може бути неоднозначним. У~нерівноважній плазмі необхідно окремо вказувати $T_e$ та $T_i$.

%% --------------------------------------------------------
\subsection{Квазінейтральність плазми}
%% --------------------------------------------------------

Ще одна принципова властивість плазми~--- \emph{квазінейтральність}. Якщо б у~середовищі була присутня велика кількість нескомпенсованих зарядів, то неможливим було б стан рівноваги, оскільки виникали б сильні поля, які приводять заряди в~рух. Тому в~плазмі густини позитивних і~негативних зарядів повинні практично збігатися за величиною. Ця властивість називається \emphx{квазінейтральністю плазми}. Відхилення від квазінейтральності приводять до виникнення сильних полів, які швидко відновлюють квазінейтральність.

На макроскопічних масштабах у~плазмі майже компенсуються позитивний і~негативний заряди:
% ~~~~~~~~~~~~~~~~~~~~~~~~~~~~~~~~~~~~~~~~~~~~~~~~~~~~~~~~
\begin{equation}
	\rho = \sum_s q_s n_s \approx 0.
\end{equation}
% ~~~~~~~~~~~~~~~~~~~~~~~~~~~~~~~~~~~~~~~~~~~~~~~~~~~~~~~~
Це не означає, що $\rho$ дорівнює нулю в~кожній точці. Невеликі порушення квазінейтральності створюють електричні поля, які швидко перерозподіляють заряд. Саме механізм цього відновлення приводить до дебаївського екранування та плазмових коливань.

Якщо концентрація іонів рівна $n_i$, а~електронів~--- $n_e$, причому кожний іон віддає в~плазму $Z$ електронів, то умова квазінейтральності записується у~вигляді $Zn_i = n_e$, або
% ~~~~~~~~~~~~~~~~~~~~~~~~~~~~~~~~~~~~~~~~~~~~~~~~~~~~~~~~
\begin{equation}\label{eq:quasineutrality}
	e n_e + (-Ze) n_i = 0.
\end{equation}
% ~~~~~~~~~~~~~~~~~~~~~~~~~~~~~~~~~~~~~~~~~~~~~~~~~~~~~~~~
Тут $e < 0$~--- заряд електрона.

Розглянемо спочатку однорідну плазму, яка в~середньому є електронейтральною. Нехай у~невеликій області виник надлишок позитивного заряду. Виникле електричне поле почне переміщувати електрони до цієї області та відштовхувати від неї позитивні іони. Отже, плазма чинить опір локальному порушенню квазінейтральності.

Характерний час, за який електрони реагують на порушення заряду, за порядком величини визначається плазмовою частотою. Характерний просторовий масштаб такого порушення визначається дебаївською довжиною. Тому для плазми важливо, щоб
% ~~~~~~~~~~~~~~~~~~~~~~~~~~~~~~~~~~~~~~~~~~~~~~~~~~~~~~~~
\begin{equation}
	\lambda_D\ll L,
\end{equation}
% ~~~~~~~~~~~~~~~~~~~~~~~~~~~~~~~~~~~~~~~~~~~~~~~~~~~~~~~~
де $L$~--- характерний макроскопічний розмір системи. Цю нерівність слід розуміти як умову того, що локальні порушення квазінейтральності швидко екрануються всередині плазми. Біля стінок, електродів та інших меж ця умова може порушуватися, і~там виникають дебаївські шари.

%% --------------------------------------------------------
\subsection{Дебаївський радіус}
%% --------------------------------------------------------

Відхилення від нейтральності може відбуватися тільки на малих відстанях і~на протязі малих проміжків часу. Оцінимо характерні просторові масштаби, на яких може відбуватися сильне розділення зарядів.

Нехай в~деякому шарі товщиною $l$ і~об'ємом $V = Sl$ відбулося розділення зарядів, причому всі електрони зібралися поблизу верхньої площини, а~іони~--- поблизу нижньої площини (рис.~\ref{pic:charge-separation}). Обмежимося випадком однократно іонізованої плазми, $Z = 1$. При цьому в~стані рівноваги $n_i = n_e \equiv n$.

\begin{figure}[h!]\centering
	% \includegraphics[width=0.5\linewidth]{charge-separation}
	\caption{Розділення зарядів у~плазмі приводить до виникнення полів, які прагнуть відновити нейтральність}
	\label{pic:charge-separation}
\end{figure}

В~результаті розділення на площинах виникають нескомпенсовані заряди з~поверхневою густиною $\sigma = neV/S = nel$. Електричне поле в~такому конденсаторі виявиться рівним $E = 4\pi\sigma = 4\pi nel$. Густина енергії електричного поля $w_E = E^2/8\pi$. Ця енергія утворюється за рахунок кінетичної енергії $w_T$ теплового руху електронів і~іонів. Маючи на увазі, що в~наведеній оцінці розглядається розділення зарядів вздовж одного напрямку (осі $x$ на рис.~\ref{pic:charge-separation}), на основі теореми про рівнорозподіл енергії за ступенями свободи маємо
% ~~~~~~~~~~~~~~~~~~~~~~~~~~~~~~~~~~~~~~~~~~~~~~~~~~~~~~~~
\begin{equation*}
	w_T = w_T^{(e)} + w_T^{(i)} = n\frac{kT}{2} + n\frac{kT}{2} = nkT.
\end{equation*}
% ~~~~~~~~~~~~~~~~~~~~~~~~~~~~~~~~~~~~~~~~~~~~~~~~~~~~~~~~

За законом збереження енергії $w_E = w_T$. Звідси випливає, що
% ~~~~~~~~~~~~~~~~~~~~~~~~~~~~~~~~~~~~~~~~~~~~~~~~~~~~~~~~
\begin{equation*}
	E = \sqrt{8\pi nkT},
\end{equation*}
% ~~~~~~~~~~~~~~~~~~~~~~~~~~~~~~~~~~~~~~~~~~~~~~~~~~~~~~~~
або
% ~~~~~~~~~~~~~~~~~~~~~~~~~~~~~~~~~~~~~~~~~~~~~~~~~~~~~~~~
\begin{equation*}
	4\pi nel = \sqrt{8\pi nkT} \quad\Rightarrow\quad l = \sqrt{\frac{kT}{2\pi ne^2}}.
\end{equation*}
% ~~~~~~~~~~~~~~~~~~~~~~~~~~~~~~~~~~~~~~~~~~~~~~~~~~~~~~~~

Величина $r_D = l/2$ або
% ~~~~~~~~~~~~~~~~~~~~~~~~~~~~~~~~~~~~~~~~~~~~~~~~~~~~~~~~
\begin{equation}\label{eq:debye-radius}
	r_D = \sqrt{\frac{kT}{8\pi ne^2}},
\end{equation}
% ~~~~~~~~~~~~~~~~~~~~~~~~~~~~~~~~~~~~~~~~~~~~~~~~~~~~~~~~
називається \emphx{дебаївським радіусом} і~вказує характерний просторовий масштаб, на якому може відбуватися суттєве відхилення від квазінейтральності.

%% --------------------------------------------------------
\subsection{Дебаївське екранування}
%% --------------------------------------------------------

Розглянемо малий пробний заряд $q$, занурений у~плазму. Для простоти припустимо, що плазма складається з~електронів та одноразово заряджених іонів, а~відхилення густин від рівноважних малі. Завдяки кулонівській взаємодії відбудеться просторовий перерозподіл частинок плазми, в~результаті чого внесений заряд буде оточений \enquote{шубою} з~частинок іншого знака заряду. Відбудеться екранування, ослаблення поля, що створюється пробним зарядом. Однак повної нейтралізації не спостерігатиметься внаслідок теплового руху частинок плазми.

У~найпростішому випадку на масштабах, характерних для електронного екранування, іони можна вважати нерухомими. Тоді електрони в електростатичному потенціалі $\varphi$ мають больцманівський розподіл
% ~~~~~~~~~~~~~~~~~~~~~~~~~~~~~~~~~~~~~~~~~~~~~~~~~~~~~~~~
\begin{equation}
	n_e=n_0\exp\left(\frac{e\varphi}{k_{\mathrm B}T_e}\right).
\end{equation}
% ~~~~~~~~~~~~~~~~~~~~~~~~~~~~~~~~~~~~~~~~~~~~~~~~~~~~~~~~
Для слабкого поля
\begin{equation*}
	\left|\frac{e\varphi}{k_{\mathrm B}T_e}\right|\ll1
\end{equation*}
одержуємо
% ~~~~~~~~~~~~~~~~~~~~~~~~~~~~~~~~~~~~~~~~~~~~~~~~~~~~~~~~
\begin{equation}
	n_e\simeq n_0
	\left(1+\frac{e\varphi}{k_{\mathrm B}T_e}\right).
\end{equation}
% ~~~~~~~~~~~~~~~~~~~~~~~~~~~~~~~~~~~~~~~~~~~~~~~~~~~~~~~~
Тоді густина заряду плазми
% ~~~~~~~~~~~~~~~~~~~~~~~~~~~~~~~~~~~~~~~~~~~~~~~~~~~~~~~~
\begin{equation}
	\rho_{\mathrm{pl}}
	=e n_0-e n_e
	\simeq-\frac{n_0e^2}{k_{\mathrm B}T_e}\varphi.
\end{equation}
% ~~~~~~~~~~~~~~~~~~~~~~~~~~~~~~~~~~~~~~~~~~~~~~~~~~~~~~~~
Разом із рівнянням Пуассона
% ~~~~~~~~~~~~~~~~~~~~~~~~~~~~~~~~~~~~~~~~~~~~~~~~~~~~~~~~
\begin{equation}
	\nabla^2\varphi=-4\pi\rho
\end{equation}
% ~~~~~~~~~~~~~~~~~~~~~~~~~~~~~~~~~~~~~~~~~~~~~~~~~~~~~~~~
це дає
% ~~~~~~~~~~~~~~~~~~~~~~~~~~~~~~~~~~~~~~~~~~~~~~~~~~~~~~~~
\begin{equation}
	\nabla^2\varphi-\frac{\varphi}{\lambda_{De}^2}
	=-4\pi q\delta(\vect r),
\end{equation}
% ~~~~~~~~~~~~~~~~~~~~~~~~~~~~~~~~~~~~~~~~~~~~~~~~~~~~~~~~
де
% ~~~~~~~~~~~~~~~~~~~~~~~~~~~~~~~~~~~~~~~~~~~~~~~~~~~~~~~~
\begin{equation}
	\tcbhighmath{\lambda_{De}
		=\sqrt{\frac{k_{\mathrm B}T_e}{4\pi n_e e^2}}}
\end{equation}
% ~~~~~~~~~~~~~~~~~~~~~~~~~~~~~~~~~~~~~~~~~~~~~~~~~~~~~~~~
називається \emph{електронною дебаївською довжиною}.

Для точкового пробного заряду розв'язок має вигляд
% ~~~~~~~~~~~~~~~~~~~~~~~~~~~~~~~~~~~~~~~~~~~~~~~~~~~~~~~~
\begin{equation}
	\varphi(r)
	=\frac{q}{r}
	\exp\left(-\frac{r}{\lambda_{De}}\right).
\end{equation}
% ~~~~~~~~~~~~~~~~~~~~~~~~~~~~~~~~~~~~~~~~~~~~~~~~~~~~~~~~
Отже, кулонівське поле зарядженої частинки в~плазмі експоненційно послаблюється на відстанях порядку дебаївської довжини. Плазма екранує локальний заряд. Цей результат і~характерне просторове зменшення потенціалу показано на \ref{fig:debye-screening}.

\begin{figure}[htbp]\centering
    \localinput{potential.tikz}
	\caption{Дебаївське екранування пробного заряду: кулонівський потенціал швидко зменшується на відстанях порядку $\lambda_D$.}\label{fig:debye-screening}
\end{figure}

Якщо електрони та іони обидва встигають перебудуватися і~мають температури $T_e$ та $T_i$, загальна дебаївська довжина визначається формулою
% ~~~~~~~~~~~~~~~~~~~~~~~~~~~~~~~~~~~~~~~~~~~~~~~~~~~~~~~~
\begin{equation}
	\tcbhighmath{
		\frac{1}{\lambda_D^2}
		=4\pi\sum_s\frac{n_s q_s^2}{k_{\mathrm B}T_s}
	}.
\end{equation}
% ~~~~~~~~~~~~~~~~~~~~~~~~~~~~~~~~~~~~~~~~~~~~~~~~~~~~~~~~
Для повністю рухомої однорідної плазми з~електронами та одноразово зарядженими іонами
% ~~~~~~~~~~~~~~~~~~~~~~~~~~~~~~~~~~~~~~~~~~~~~~~~~~~~~~~~
\begin{equation}
	\frac{1}{\lambda_D^2}
	=4\pi\left(
	\frac{n_e e^2}{k_{\mathrm B}T_e}
	+\frac{n_i e^2}{k_{\mathrm B}T_i}\right).
\end{equation}
% ~~~~~~~~~~~~~~~~~~~~~~~~~~~~~~~~~~~~~~~~~~~~~~~~~~~~~~~~
При $n_e=n_i=n$ і~$T_e=T_i=T$
% ~~~~~~~~~~~~~~~~~~~~~~~~~~~~~~~~~~~~~~~~~~~~~~~~~~~~~~~~
\begin{equation}
	\lambda_D
	=\sqrt{\frac{k_{\mathrm B}T}{8\pi ne^2}}.
\end{equation}
% ~~~~~~~~~~~~~~~~~~~~~~~~~~~~~~~~~~~~~~~~~~~~~~~~~~~~~~~~

У~літературі з~фізики плазми під \enquote{дебаївською довжиною} дуже часто мають на увазі саме електронну величину $\lambda_{De}$. У~задачах потрібно тому завжди перевіряти, які компоненти плазми встигають перебудуватися на розглядуваному часовому масштабі. Таким чином, на відстанях порядку дебаївського радіуса відбувається сильне екранування внесеного в~плазму заряду. Тому величину $\lambda_D$ називають також \emphx{радіусом екранування}.

%% --------------------------------------------------------
\subsection{Дебаївська сфера і~умова колективності}
%% --------------------------------------------------------

\emphx{Дебаївською сферою} називають область в~плазмі, обмежену сферою з~радіусом, рівним дебаївському радіусу $r_D$. Якщо концентрація частинок в~плазмі є $n$, то в~дебаївській сфері виявиться
% ~~~~~~~~~~~~~~~~~~~~~~~~~~~~~~~~~~~~~~~~~~~~~~~~~~~~~~~~
\begin{equation*}
	N \sim \frac{4}{3}\pi r_D^3 n
\end{equation*}
% ~~~~~~~~~~~~~~~~~~~~~~~~~~~~~~~~~~~~~~~~~~~~~~~~~~~~~~~~
частинок. В~плазмі з~концентрацією частинок $n$ середня відстань між ними становить $d \sim n^{-1/3}$. Відповідно маємо
% ~~~~~~~~~~~~~~~~~~~~~~~~~~~~~~~~~~~~~~~~~~~~~~~~~~~~~~~~
\begin{equation*}
	N \sim (r_D/d)^3.
\end{equation*}
% ~~~~~~~~~~~~~~~~~~~~~~~~~~~~~~~~~~~~~~~~~~~~~~~~~~~~~~~~

Число частинок у~сфері радіуса $\lambda_{De}$ оцінюється як
% ~~~~~~~~~~~~~~~~~~~~~~~~~~~~~~~~~~~~~~~~~~~~~~~~~~~~~~~~
\begin{equation}
	N_D=\frac{4\pi}{3}n_e\lambda_{De}^3.
\end{equation}
% ~~~~~~~~~~~~~~~~~~~~~~~~~~~~~~~~~~~~~~~~~~~~~~~~~~~~~~~~
Для звичайного плазмового наближення потрібно
% ~~~~~~~~~~~~~~~~~~~~~~~~~~~~~~~~~~~~~~~~~~~~~~~~~~~~~~~~
\begin{equation}
	\tcbhighmath{N_D\gg1.}
\end{equation}
% ~~~~~~~~~~~~~~~~~~~~~~~~~~~~~~~~~~~~~~~~~~~~~~~~~~~~~~~~
Це означає, що екранування створюється не окремою частинкою, а статистично усередненим полем великої кількості частинок.

Нехай $a$~--- середня відстань між частинками,
% ~~~~~~~~~~~~~~~~~~~~~~~~~~~~~~~~~~~~~~~~~~~~~~~~~~~~~~~~
\begin{equation}
	a\sim n^{-1/3}.
\end{equation}
% ~~~~~~~~~~~~~~~~~~~~~~~~~~~~~~~~~~~~~~~~~~~~~~~~~~~~~~~~
Тоді умова $N_D\gg1$ означає
% ~~~~~~~~~~~~~~~~~~~~~~~~~~~~~~~~~~~~~~~~~~~~~~~~~~~~~~~~
\begin{equation*}
	\lambda_D\gg a.
\end{equation*}
% ~~~~~~~~~~~~~~~~~~~~~~~~~~~~~~~~~~~~~~~~~~~~~~~~~~~~~~~~

На \ref{fig:debye-sphere} схематично показано дебаївську сферу: саме велика кількість частинок у~ній робить екранування колективним.

\begin{SCfigure}[2][h!]\centering
	\localinput{debye_sphere.tikz}
	\caption{Дебаївська сфера та число частинок у~ній. У~колективному режимі $N_D\gg1$: екранування створюється багатьма частинками.}\label{fig:debye-sphere}
\end{SCfigure}

Отже, в~плазмі виникає принципово важлива ієрархія масштабів: мікроскопічна відстань між частинками набагато менша за масштаб колективного екранування, а~той, у~свою чергу, має бути малим порівняно з~розміром плазми:
% ~~~~~~~~~~~~~~~~~~~~~~~~~~~~~~~~~~~~~~~~~~~~~~~~~~~~~~~~
\begin{equation}
	a\ll\lambda_D\ll L.
\end{equation}
% ~~~~~~~~~~~~~~~~~~~~~~~~~~~~~~~~~~~~~~~~~~~~~~~~~~~~~~~~

Ця нерівність є одним із найкорисніших критеріїв для інтуїтивного розуміння плазми. Вона одночасно означає наявність великої кількості частинок у~дебаївській сфері та можливість локально відновлювати квазінейтральність.

Виділяють два граничних випадки.

1) $N \gg 1$. Це значить, що $r_D \gg d \sim n^{-1/3}$. Іншими словами, в~дебаївській сфері знаходиться багато частинок. Оскільки екранування на відстанях, менших $r_D$, відносно невелике, то домінують колективні ефекти, зумовлені далекодіючими кулонівськими силами.

Перепишемо нерівність $r_D \gg d$, врахувавши визначення $r_D$:
% ~~~~~~~~~~~~~~~~~~~~~~~~~~~~~~~~~~~~~~~~~~~~~~~~~~~~~~~~
\begin{equation*}
	r_D \sim \sqrt{\frac{kT}{8\pi ne^2}} \gg n^{-1/3} \quad\Rightarrow\quad n \ll \left(\frac{kT}{e^2}\right)^3.
\end{equation*}
% ~~~~~~~~~~~~~~~~~~~~~~~~~~~~~~~~~~~~~~~~~~~~~~~~~~~~~~~~

Таким чином, нерівність $N \gg 1$ означає, що густина плазми достатньо мала. Плазма, яка задовольняє сформульовану умову, називається \emphx{розрідженою}.

Заметимо, що умову розрідженої плазми можна переписати у~вигляді
% ~~~~~~~~~~~~~~~~~~~~~~~~~~~~~~~~~~~~~~~~~~~~~~~~~~~~~~~~
\begin{equation*}
	e^2/d \ll kT, \quad d \sim n^{-1/3}.
\end{equation*}
% ~~~~~~~~~~~~~~~~~~~~~~~~~~~~~~~~~~~~~~~~~~~~~~~~~~~~~~~~

Смисл цієї умови в~тому, що потенціальна енергія взаємодії частинок мала в~порівнянні з~їхньою кінетичною енергією. В~даному випадку плазма подібна ідеальному газу. Тому таку плазму називають також \emphx{ідеальною}.

2) $N \ll 1$. Відповідно $r_D \ll d \sim n^{-1/3}$. В~цьому випадку густина плазми велика:
% ~~~~~~~~~~~~~~~~~~~~~~~~~~~~~~~~~~~~~~~~~~~~~~~~~~~~~~~~
\begin{equation*}
	n \gg (kT/e^2)^3.
\end{equation*}
% ~~~~~~~~~~~~~~~~~~~~~~~~~~~~~~~~~~~~~~~~~~~~~~~~~~~~~~~~

Така плазма називається \emphx{густою}. В~ній просторове розділення заряду можливо тільки на відстані порядку середньої відстані між частинками. Але на таких відстанях квазінейтральність і~так порушується. Густа плазма не може розглядатися як ідеальний газ.

%% --------------------------------------------------------
\subsection{Слабкозв'язана та сильнозв'язана плазма}
%% --------------------------------------------------------

Для оцінки важливості парної кулонівської взаємодії введемо параметр зв'язку
% ~~~~~~~~~~~~~~~~~~~~~~~~~~~~~~~~~~~~~~~~~~~~~~~~~~~~~~~~
\begin{equation}
	\Gamma=
	\frac{e^2}{a\,k_{\mathrm B}T}.
\end{equation}
% ~~~~~~~~~~~~~~~~~~~~~~~~~~~~~~~~~~~~~~~~~~~~~~~~~~~~~~~~
Якщо
% ~~~~~~~~~~~~~~~~~~~~~~~~~~~~~~~~~~~~~~~~~~~~~~~~~~~~~~~~
\begin{equation}
	\Gamma\ll1,
\end{equation}
% ~~~~~~~~~~~~~~~~~~~~~~~~~~~~~~~~~~~~~~~~~~~~~~~~~~~~~~~~
середня енергія кулонівської взаємодії мала порівняно з~тепловою енергією. Таку плазму називають \emph{слабкозв'язаною} або \emph{слабкокорельованою}. У~цьому випадку її термодинамічні властивості в~першому наближенні близькі до властивостей ідеального газу:
% ~~~~~~~~~~~~~~~~~~~~~~~~~~~~~~~~~~~~~~~~~~~~~~~~~~~~~~~~
\begin{equation}
	p\simeq \sum_s n_s k_{\mathrm B}T_s.
\end{equation}
% ~~~~~~~~~~~~~~~~~~~~~~~~~~~~~~~~~~~~~~~~~~~~~~~~~~~~~~~~
Водночас це зовсім не означає відсутності колективних явищ: саме слабка парна взаємодія разом із далекодійністю кулонівського поля породжує сильні колективні ефекти.

При $\Gamma\gtrsim1$ кулонівська кореляція стає істотною, і ідеально-газове наближення вже недостатнє. Така плазма належить до сильнозв'язаного режиму. За ще більшої густини або меншої температури важливими можуть стати квантові ефекти.

%% --------------------------------------------------------
\subsection{Амбіполярна дифузія}
%% --------------------------------------------------------

Якщо в плазмі є градієнт концентрації, то, незважаючи на різницю коефіцієнтів дифузії, позитивні й негативні іони не дифундують з різними швидкостями, а завдяки квазінейтральності плазми переміщаються з однією й тією ж середньою швидкістю в одному й тому ж напрямку. Така дифузія називається \emphx{амбіполярною}. Фізична причина амбіполярного характеру дифузії плазми полягає в тому, що через різницю в коефіцієнтах дифузії та рухливості іонів у плазмі виникає деяке розділення електричних зарядів і пов'язане з ним електричне поле, яке сповільнює дифузію іонів одного знака й прискорює дифузію іонів протилежного знака. У результаті швидкості дифузії позитивних і негативних іонів вирівнюються. Відновлення порушеної квазінейтральності плазми аналогічне появі відновлювальних сил у пружних тілах при їхніх деформаціях. З цим пов'язана можливість різноманітних \emphx{колективних коливань} плазми, значно більш різноманітних, ніж у газах, що складаються з нейтральних частинок.

%% --------------------------------------------------------
\section{Плазмові коливання та плазмова частота}\label{sec:plasma-oscillations}
%% --------------------------------------------------------

Дебаївське екранування описує просторову відповідь плазми. Тепер розглянемо її часову відповідь. Нехай у~шарі плазми товщиною $x$ відбулося розділення зарядів: всі електрони зібралися поблизу верхньої площини, а~іони~--- поблизу нижньої площини (рис.~\ref{pic:plasma-oscillations}). В~сформованому конденсаторі об'ємом $V = Sx$ виникло електричне поле ($e < 0$):
% ~~~~~~~~~~~~~~~~~~~~~~~~~~~~~~~~~~~~~~~~~~~~~~~~~~~~~~~~
\begin{equation*}
	E = 4\pi\sigma, \quad \sigma = -\frac{neV}{S} = -nex \quad\Rightarrow\quad E = -4\pi nex.
\end{equation*}
% ~~~~~~~~~~~~~~~~~~~~~~~~~~~~~~~~~~~~~~~~~~~~~~~~~~~~~~~~

Тут $\sigma$~--- поверхнева густина заряду позитивно зарядженої (нижньої) пластини конденсатора. Це поле повідомляє кожному електрону прискорення $a = eE/m$. Для випадку, коли електрони знаходяться у~верхньому шарі (як показано на рис.~\ref{pic:plasma-oscillations}), їхнє прискорення напрямлене вниз: $E > 0$, $a < 0$.

\begin{figure}[h!]\centering
	% \includegraphics[width=0.5\linewidth]{plasma-oscillations}
	\caption{До розрахунку частоти плазмових коливань}
	\label{pic:plasma-oscillations}
\end{figure}

Запишемо рівняння руху електронів:
% ~~~~~~~~~~~~~~~~~~~~~~~~~~~~~~~~~~~~~~~~~~~~~~~~~~~~~~~~
\begin{equation}\label{eq:plasma-oscillation-eq}
	m\ddot{x} = eE \quad\Rightarrow\quad m\frac{d^2x}{dt^2} = -4\pi ne^2 x \quad\Rightarrow\quad \frac{d^2x}{dt^2} = -\omega_p^2 x.
\end{equation}
% ~~~~~~~~~~~~~~~~~~~~~~~~~~~~~~~~~~~~~~~~~~~~~~~~~~~~~~~~

Тут введено позначення
% ~~~~~~~~~~~~~~~~~~~~~~~~~~~~~~~~~~~~~~~~~~~~~~~~~~~~~~~~
\begin{equation}\label{eq:plasma-frequency}
	\omega_p = \sqrt{4\pi ne^2/m}.
\end{equation}
% ~~~~~~~~~~~~~~~~~~~~~~~~~~~~~~~~~~~~~~~~~~~~~~~~~~~~~~~~

Ця величина називається \emphx{плазмовою} (або \emphx{ленгмюрівською}\footnote{Назва походить від імені І.~Ленгмюра, який досліджував газовий розряд і~запропонував термін \enquote{плазма}.}) \emphx{частотою}.

Як видно з~отриманого рівняння, електрони здійснюють гармонічні коливання відносно іонів з~плазмовою частотою. Ці коливання називають \emphx{плазмовими}. Їх амплітуда порядку дебаївського радіуса. Останнє пов'язано з~тим, що величина $r_D$ є характерний масштаб просторового розділення зарядів у~плазмі (див.~\eqref{eq:debye-radius}).

Той самий результат можна отримати без уявного конденсатора. З~рівняння неперервності
% ~~~~~~~~~~~~~~~~~~~~~~~~~~~~~~~~~~~~~~~~~~~~~~~~~~~~~~~~
\begin{equation}
	\frac{\partial\rho}{\partial t}
	+\nabla\!\cdot\!\vect j=0
\end{equation}
% ~~~~~~~~~~~~~~~~~~~~~~~~~~~~~~~~~~~~~~~~~~~~~~~~~~~~~~~~
та рівняння руху електронів
% ~~~~~~~~~~~~~~~~~~~~~~~~~~~~~~~~~~~~~~~~~~~~~~~~~~~~~~~~
\begin{equation}
	m_e\frac{\partial\vect v}{\partial t}=-e\vect E
\end{equation}
% ~~~~~~~~~~~~~~~~~~~~~~~~~~~~~~~~~~~~~~~~~~~~~~~~~~~~~~~~
за умови сталих $n_e$ та нехтування магнітним полем одержуємо
% ~~~~~~~~~~~~~~~~~~~~~~~~~~~~~~~~~~~~~~~~~~~~~~~~~~~~~~~~
\begin{equation}
	\frac{\partial^2\rho}{\partial t^2}
	+\omega_{pe}^2\rho=0.
\end{equation}
% ~~~~~~~~~~~~~~~~~~~~~~~~~~~~~~~~~~~~~~~~~~~~~~~~~~~~~~~~
Таким чином, плазмова частота є власною частотою електронного компонента плазми. Фізичний механізм цього коливання — зміщення електронного компонента відносно позитивного фону — схематично показано на \ref{fig:plasma-oscillation}.

\begin{figure}[htbp]\centering
	%---------------------------------------------------------
	\begin{subfigure}[t]{0.5\linewidth}\centering
		\localinput{plasma_ions.tikz}
	\end{subfigure}%
	%---------------------------------------------------------
	\begin{subfigure}[t]{0.5\linewidth}\centering
		\localinput{condensator.tikz}
	\end{subfigure}
	%---------------------------------------------------------
	\caption{Модель плазмових коливань: електронний компонент зміщений відносно нерухомого позитивного фону. Виникає повертаюче електричне поле.}\label{fig:plasma-oscillation}
\end{figure}

Якщо врахувати рух іонів, то для холодної двокомпонентної плазми частота довгохвильових поздовжніх коливань визначається
% ~~~~~~~~~~~~~~~~~~~~~~~~~~~~~~~~~~~~~~~~~~~~~~~~~~~~~~~~
\begin{equation}
	\omega_p^2
	=4\pi e^2\left(
	\frac{n_e}{m_e}+\frac{Z^2n_i}{m_i}\right).
\end{equation}
% ~~~~~~~~~~~~~~~~~~~~~~~~~~~~~~~~~~~~~~~~~~~~~~~~~~~~~~~~
Для $m_i\gg m_e$ іонний внесок малий, і
\begin{equation*}
	\omega_p\simeq\omega_{pe}.
\end{equation*}

У~реальній плазмі коливання можуть загасати внаслідок зіткнень та інших механізмів. Крім того, якщо врахувати просторову неоднорідність і~тепловий рух електронів, частота залежить від хвильового числа. Тому величина $\omega_{pe}$ є насамперед базовою масштабною частотою плазми, а~не універсальною частотою будь-якої плазмової хвилі.

%% --------------------------------------------------------
\subsection{Зв'язок дебаївської довжини з~плазмовою частотою}
%% --------------------------------------------------------

Дві найважливіші характеристики електронної плазми пов'язані простою формулою
% ~~~~~~~~~~~~~~~~~~~~~~~~~~~~~~~~~~~~~~~~~~~~~~~~~~~~~~~~
\begin{equation}
	\lambda_{De}
	=\frac{v_{Te}}{\omega_{pe}},
	\qquad
	v_{Te}=\sqrt{\frac{k_{\mathrm B}T_e}{m_e}},
\end{equation}
% ~~~~~~~~~~~~~~~~~~~~~~~~~~~~~~~~~~~~~~~~~~~~~~~~~~~~~~~~
де $v_{Te}$~--- характерна теплова швидкість електронів.

Отже, дебаївська довжина~--- це відстань, яку електрон із характерною тепловою швидкістю проходить за один час порядку $\omega_{pe}^{-1}$. Це дає корисну фізичну інтерпретацію двох фундаментальних масштабів:
\begin{equation*}
	\tcbhighmath{\text{час колективної реакції}\sim\omega_{pe}^{-1}},
	\qquad
	\tcbhighmath{\text{довжина колективної реакції}\sim\lambda_{De}}.
\end{equation*}

%% --------------------------------------------------------
\section{Діелектрична проникність плазми}\label{sec:plasma-permittivity}\hypertarget{plasma-permittivity}{}
%% --------------------------------------------------------

Плазма реагує не лише на статичне поле, а~й~на змінне електромагнітне поле. Розглянемо, як плазма відповідає на зовнішню електромагнітну хвилю. Це дасть змогу зрозуміти, за яких умов плазма є прозорою, а~за яких~--- відбиває випромінювання.

%% --------------------------------------------------------
\subsection{Діелектрична проникність плазми}
%% --------------------------------------------------------

Нехай на плазму падає випромінювання~--- електромагнітна хвиля. Плазма являє собою квазінейтральну суміш електронів і~іонів, причому маса іона приблизно в~$10^4$ разів більша маси електрона.

Тому в~полі зовнішньої електромагнітної хвилі рухаються головним чином електрони. Іони ж можна вважати такими, що покояться (струм, який вони створюють, виявляється малим порівняно зі струмом, який створюється електронами). У~типових умовах швидкість руху електронів мала ($v \ll c$), так що можна знехтувати силою Лоренца, яка виникає з~боку магнітного поля хвилі. Нарешті, вважаючи амплітуду коливань електрона малою в~порівнянні з~довжиною хвилі випромінювання, можна розглядати рух електрона як рух в~однорідному електричному полі хвилі, періодично змінному з~часом:
% ~~~~~~~~~~~~~~~~~~~~~~~~~~~~~~~~~~~~~~~~~~~~~~~~~~~~~~~~
\begin{equation*}
	E = E_0 \cos\omega t.
\end{equation*}
% ~~~~~~~~~~~~~~~~~~~~~~~~~~~~~~~~~~~~~~~~~~~~~~~~~~~~~~~~

Виберемо систему координат так, щоб вісь $x$ була напрямлена вздовж вектора електричного поля. В~цих умовах рух окремого електрона описується рівнянням
% ~~~~~~~~~~~~~~~~~~~~~~~~~~~~~~~~~~~~~~~~~~~~~~~~~~~~~~~~
\begin{equation*}
	m\ddot{x} = eE = eE_0\cos\omega t,
\end{equation*}
% ~~~~~~~~~~~~~~~~~~~~~~~~~~~~~~~~~~~~~~~~~~~~~~~~~~~~~~~~
розв'язок якого має вигляд
% ~~~~~~~~~~~~~~~~~~~~~~~~~~~~~~~~~~~~~~~~~~~~~~~~~~~~~~~~
\begin{equation*}
	x = -\frac{eE_0}{m\omega^2}\cos\omega t = -\frac{eE}{m\omega^2}.
\end{equation*}
% ~~~~~~~~~~~~~~~~~~~~~~~~~~~~~~~~~~~~~~~~~~~~~~~~~~~~~~~~

У~розглядуваній ситуації електрони плазми можна уподібнити зв'язаним зарядам у~діелектрику, оскільки вони не переміщаються на більші відстані, а~лише здійснюють коливання відносно свого вихідного положення. Тому зміщення, індуковане електромагнітною хвилею, приводить до появи ефективного дипольного моменту електрона
% ~~~~~~~~~~~~~~~~~~~~~~~~~~~~~~~~~~~~~~~~~~~~~~~~~~~~~~~~
\begin{equation*}
	p = ex = -e^2E/m\omega^2,
\end{equation*}
% ~~~~~~~~~~~~~~~~~~~~~~~~~~~~~~~~~~~~~~~~~~~~~~~~~~~~~~~~
або, у~векторному вигляді,
% ~~~~~~~~~~~~~~~~~~~~~~~~~~~~~~~~~~~~~~~~~~~~~~~~~~~~~~~~
\begin{equation*}
	\vect{p} = -e^2\vect{E}/m\omega^2.
\end{equation*}
% ~~~~~~~~~~~~~~~~~~~~~~~~~~~~~~~~~~~~~~~~~~~~~~~~~~~~~~~~

Якщо в~одиниці об'єму плазми мається $N$ електронів, то вектор поляризації (дипольний момент одиниці об'єму) виявиться рівним
% ~~~~~~~~~~~~~~~~~~~~~~~~~~~~~~~~~~~~~~~~~~~~~~~~~~~~~~~~
\begin{equation*}
	\vect{P} = N\vect{p} = -\frac{Ne^2}{m\omega^2}\vect{E} \equiv \alpha\vect{E}.
\end{equation*}
% ~~~~~~~~~~~~~~~~~~~~~~~~~~~~~~~~~~~~~~~~~~~~~~~~~~~~~~~~

Маючи на увазі, що $\vect{D} = \vect{E} + 4\pi\vect{P} = \varepsilon\vect{E}$, заключаємо, що діелектрична проникність плазми рівна
% ~~~~~~~~~~~~~~~~~~~~~~~~~~~~~~~~~~~~~~~~~~~~~~~~~~~~~~~~
\begin{equation}\label{eq:plasma-permittivity}
	\varepsilon = 1 + 4\pi\alpha = 1 - \frac{4\pi Ne^2}{m\omega^2}.
\end{equation}
% ~~~~~~~~~~~~~~~~~~~~~~~~~~~~~~~~~~~~~~~~~~~~~~~~~~~~~~~~

Оскільки $\sqrt{4\pi Ne^2/m} = \omega_p$~--- плазмова частота, то остаточно отримуємо
% ~~~~~~~~~~~~~~~~~~~~~~~~~~~~~~~~~~~~~~~~~~~~~~~~~~~~~~~~
\begin{equation}\label{eq:plasma-permittivity-final}
	\tcbhighmath{\varepsilon_r(\omega) = 1 - \omega_p^2/\omega^2}.
\end{equation}
% ~~~~~~~~~~~~~~~~~~~~~~~~~~~~~~~~~~~~~~~~~~~~~~~~~~~~~~~~

Таким чином, на відміну від випадку діелектриків в~електростатичному полі, діелектрична проникність плазми виявляється менше одиниці: $\varepsilon < 1$. Більше того, вона від'ємна при частотах $\omega < \omega_p$.

Той самий результат можна отримати, якщо записати електронний струм
% ~~~~~~~~~~~~~~~~~~~~~~~~~~~~~~~~~~~~~~~~~~~~~~~~~~~~~~~~
\begin{equation}
	m_e\frac{d\vect v}{dt}=-e\vect E,
	\qquad
	\vect j=-n_e e\vect v.
\end{equation}
% ~~~~~~~~~~~~~~~~~~~~~~~~~~~~~~~~~~~~~~~~~~~~~~~~~~~~~~~~
Для гармонічного поля $\vect E\propto e^{-i\omega t}$ маємо
% ~~~~~~~~~~~~~~~~~~~~~~~~~~~~~~~~~~~~~~~~~~~~~~~~~~~~~~~~
\begin{equation}
	\vect j
	=i\frac{n_e e^2}{m_e\omega}\vect E.
\end{equation}
% ~~~~~~~~~~~~~~~~~~~~~~~~~~~~~~~~~~~~~~~~~~~~~~~~~~~~~~~~
Отже, плазма має частотно залежну електропровідність
% ~~~~~~~~~~~~~~~~~~~~~~~~~~~~~~~~~~~~~~~~~~~~~~~~~~~~~~~~
\begin{equation}
	\sigma(\omega)
	=i\frac{\omega_{pe}^2}{4\pi\omega}.
\end{equation}
% ~~~~~~~~~~~~~~~~~~~~~~~~~~~~~~~~~~~~~~~~~~~~~~~~~~~~~~~~

%% --------------------------------------------------------
\subsection{Відбиття електромагнітних хвиль від плазми}
%% --------------------------------------------------------

Як випливає з~рівнянь Максвелла, для плоских хвиль
% ~~~~~~~~~~~~~~~~~~~~~~~~~~~~~~~~~~~~~~~~~~~~~~~~~~~~~~~~
\begin{equation}\label{eq:plane-wave}
	E = E_0 \exp(ikx - i\omega t)
\end{equation}
% ~~~~~~~~~~~~~~~~~~~~~~~~~~~~~~~~~~~~~~~~~~~~~~~~~~~~~~~~
частота $\omega$ і~хвильове число $k$ зв'язані співвідношенням
% ~~~~~~~~~~~~~~~~~~~~~~~~~~~~~~~~~~~~~~~~~~~~~~~~~~~~~~~~
\begin{equation*}
	k^2 = \frac{\omega^2}{c^2}\varepsilon(\omega).
\end{equation*}
% ~~~~~~~~~~~~~~~~~~~~~~~~~~~~~~~~~~~~~~~~~~~~~~~~~~~~~~~~

Підставляючи сюди вираз~\eqref{eq:plasma-permittivity-final} для діелектричної проникності, отримуємо
% ~~~~~~~~~~~~~~~~~~~~~~~~~~~~~~~~~~~~~~~~~~~~~~~~~~~~~~~~
\begin{equation*}
	k^2 = \frac{\omega^2}{c^2}\left(1 - \frac{\omega_p^2}{\omega^2}\right).
\end{equation*}
% ~~~~~~~~~~~~~~~~~~~~~~~~~~~~~~~~~~~~~~~~~~~~~~~~~~~~~~~~

Останнє співвідношення можна переписати у~вигляді
% ~~~~~~~~~~~~~~~~~~~~~~~~~~~~~~~~~~~~~~~~~~~~~~~~~~~~~~~~
\begin{equation}\label{eq:plasma-dispersion}
	\tcbhighmath{\omega^2 = \omega_p^2 + c^2k^2}.
\end{equation}
% ~~~~~~~~~~~~~~~~~~~~~~~~~~~~~~~~~~~~~~~~~~~~~~~~~~~~~~~~

Цей закон дисперсії показано на рис.~\ref{pic:plasma-dispersion}.

Із отриманого співвідношення випливає, зокрема, що довжина хвилі рівна
% ~~~~~~~~~~~~~~~~~~~~~~~~~~~~~~~~~~~~~~~~~~~~~~~~~~~~~~~~
\begin{equation*}
	\lambda = \frac{2\pi}{k} = \frac{2\pi c}{\sqrt{\omega^2 - \omega_p^2}}.
\end{equation*}
% ~~~~~~~~~~~~~~~~~~~~~~~~~~~~~~~~~~~~~~~~~~~~~~~~~~~~~~~~

Нехай на шар плазми падає хвиля з~частотою $\omega < \omega_p$ (рис.~\ref{pic:plasma-reflection}). В~цьому випадку
% ~~~~~~~~~~~~~~~~~~~~~~~~~~~~~~~~~~~~~~~~~~~~~~~~~~~~~~~~
\begin{equation*}
	k^2 = \frac{\omega^2 - \omega_p^2}{c^2} < 0.
\end{equation*}
% ~~~~~~~~~~~~~~~~~~~~~~~~~~~~~~~~~~~~~~~~~~~~~~~~~~~~~~~~

Покладемо $k = i\beta$. Тоді вираз для хвилі~\eqref{eq:plane-wave} прийме вигляд
% ~~~~~~~~~~~~~~~~~~~~~~~~~~~~~~~~~~~~~~~~~~~~~~~~~~~~~~~~
\begin{equation*}
	E = E_0 \exp(-\beta x - i\omega t).
\end{equation*}
% ~~~~~~~~~~~~~~~~~~~~~~~~~~~~~~~~~~~~~~~~~~~~~~~~~~~~~~~~

\begin{figure}[h!]\centering
	% \includegraphics[width=0.6\linewidth]{plasma-dispersion}
	\caption{Закон дисперсії для плазми}
	\label{pic:plasma-dispersion}
\end{figure}

\begin{figure}[h!]\centering
	% \includegraphics[width=0.6\linewidth]{plasma-reflection}
	\caption{Відбиття хвилі від шару плазми}
	\label{pic:plasma-reflection}
\end{figure}

Це значить, що амплітуда хвилі експоненціально убуває~--- хвиля затухає в~міру заглиблення в~плазму.

\begin{figure}[h!]\centering
	% \includegraphics[width=0.5\linewidth]{dispersion-law}
	\caption{Закон дисперсії хвиль у~плазмі}
	\label{pic:dispersion-law}
\end{figure}

\begin{figure}[h!]\centering
	% \includegraphics[width=0.6\linewidth]{wave-incidence}
	\caption{Електромагнітна хвиля, що падає з~вакууму на шар плазми, яка займає область $x > 0$}
	\label{pic:wave-incidence}
\end{figure}

Таким чином, у~плазмі можуть розповсюджуватися тільки хвилі з~частотою $\omega > \omega_p$. Для хвиль же з~частотою $\omega < \omega_p$ плазма грає роль зеркала.

Якщо в~плазмі енергія хвилі поглинається слабо, то майже вся хвиля відіб'ється назад. Дійсно, параметри відбитої хвилі (амплітуду і~фазу) можна знайти за допомогою формул Френеля. Для випадку нормального падіння амплітудний коефіцієнт відбиття від шару плазми нескінченної товщини рівен
% ~~~~~~~~~~~~~~~~~~~~~~~~~~~~~~~~~~~~~~~~~~~~~~~~~~~~~~~~
\begin{equation*}
	r = \frac{1 - \sqrt{\varepsilon}}{1 + \sqrt{\varepsilon}}, \quad \varepsilon = 1 - \frac{\omega_p^2}{\omega^2}.
\end{equation*}
% ~~~~~~~~~~~~~~~~~~~~~~~~~~~~~~~~~~~~~~~~~~~~~~~~~~~~~~~~

При $\omega < \omega_p$ потрібно положити $\sqrt{\varepsilon} = i\kappa$. Тоді коефіцієнт відбиття за інтенсивністю дається виразом
% ~~~~~~~~~~~~~~~~~~~~~~~~~~~~~~~~~~~~~~~~~~~~~~~~~~~~~~~~
\begin{equation*}
	R = |r|^2 = \left|\frac{1 - i\kappa}{1 + i\kappa}\right|^2 = 1.
\end{equation*}
% ~~~~~~~~~~~~~~~~~~~~~~~~~~~~~~~~~~~~~~~~~~~~~~~~~~~~~~~~

Якщо ж у~плазмі суттєву роль починають грати джоулеві втрати, то коефіцієнт відбиття виявляється менше одиниці.

Якщо електромагнітна хвиля падає на шар плазми під деяким кутом до нормалі $\theta$, то відбиття може відбутися і~при частотах більших, ніж $\omega_p$. Дійсно, при $\omega > \omega_p$ плазма являє собою середовище з~дійсним показником заломлення $n = \sqrt{\varepsilon} < 1$. Тоді умова повного відбиття є
% ~~~~~~~~~~~~~~~~~~~~~~~~~~~~~~~~~~~~~~~~~~~~~~~~~~~~~~~~
\begin{equation*}
	\theta \ge \theta_{\text{в.о.}}, \quad \sin\theta_{\text{в.о.}} = n.
\end{equation*}
% ~~~~~~~~~~~~~~~~~~~~~~~~~~~~~~~~~~~~~~~~~~~~~~~~~~~~~~~~

Це~--- ефект повного внутрішнього відбиття, оскільки вакуум являється середовищем оптично більш густим, ніж плазма:
% ~~~~~~~~~~~~~~~~~~~~~~~~~~~~~~~~~~~~~~~~~~~~~~~~~~~~~~~~
\begin{equation*}
	n_{\text{вак}} = 1 > n_{\text{плазма}} = \sqrt{\varepsilon}.
\end{equation*}
% ~~~~~~~~~~~~~~~~~~~~~~~~~~~~~~~~~~~~~~~~~~~~~~~~~~~~~~~~

В~границі $\omega \to \omega_p$, коли ще $\omega < \omega_p$, кут повного внутрішнього відбиття обертається в~нуль, і~плазма відбиває хвилі при будь-яких кутах падіння.

Нехай концентрація електронів повільно росте в~напрямку осі $x$, як показано на рис.~\ref{pic:gradient-plasma}. Як правило, концентрація змінюється незначно на відстанях порядку довжини хвилі випромінювання. Тому можна приблизно визначити локальне значення хвильового числа формулою
% ~~~~~~~~~~~~~~~~~~~~~~~~~~~~~~~~~~~~~~~~~~~~~~~~~~~~~~~~
\begin{equation*}
	k^2(x) = \frac{1}{c^2}\left[\omega^2 - \frac{4\pi e^2}{m}N(x)\right].
\end{equation*}
% ~~~~~~~~~~~~~~~~~~~~~~~~~~~~~~~~~~~~~~~~~~~~~~~~~~~~~~~~

\begin{figure}[h!]\centering
	% \includegraphics[width=0.6\linewidth]{gradient-plasma}
	\caption{Електромагнітна хвиля падає на плазму зі зростаючою концентрацією електронів}
	\label{pic:gradient-plasma}
\end{figure}

Для заданої частоти хвилі хвильове число обертається в~нуль у~точці, де концентрація електронів досягає критичного значення
% ~~~~~~~~~~~~~~~~~~~~~~~~~~~~~~~~~~~~~~~~~~~~~~~~~~~~~~~~
\begin{equation*}
	N_{\text{кр}} = m\omega^2/4\pi e^2.
\end{equation*}
% ~~~~~~~~~~~~~~~~~~~~~~~~~~~~~~~~~~~~~~~~~~~~~~~~~~~~~~~~

Це співвідношення визначає ефективне положення \enquote{плазмового зеркала}.

Ефект відбиття хвиль від плазми використовується в~дальній радіозв'язку. Іоносфера (шар атмосфери, що починається на висоті приблизно 50~км від поверхні Землі) являє собою шар плазми, який формується завдяки іонізуючій дії ультрафіолетової і~рентгенівської сонячної радіації, а~також космічного випромінювання. Хоча ступінь іонізації газу в~іоносфері невисока, концентрація електронів (до $10^3$--$10^5$ електронів в~1~см$^3$) достатньо велика, щоб виявилося можливим відбиття довгих ($\lambda \sim$ декілька км) радіохвиль. Це дозволяє передавати радіосигнали на приймачі, які знаходяться поза зоною прямої видимості, як проілюстровано якісно на рис.~\ref{pic:radio-communication}.

\begin{figure}[h!]\centering
	% \includegraphics[width=0.5\linewidth]{radio-communication}
	\caption{Дальній радіозв'язок завдяки відбиттю радіохвиль від іоносфери}
	\label{pic:radio-communication}
\end{figure}

Це явище є прямим проявом колективної реакції електронів і~пов'язує поняття плазмової частоти з~електромагнітною оптикою. Вигляд дисперсійної залежності та гранична частота поширення наведені на \ref{fig:plasma-dispersion}.

\begin{figure}[htbp]\centering
    \localinput{dispersion.tikz}
	\caption{Дисперсія поперечної електромагнітної хвилі в~холодній беззіткнювальній плазмі. Частота $\omega_{pe}$ є границею поширення в~однорідній плазмі.}\label{fig:plasma-dispersion}
\end{figure}

%% --------------------------------------------------------
\section{Провідність плазми}\label{sec:plasma-conductivity}
%% --------------------------------------------------------

Найрухливішими носіями струму в~плазмі є електрони. У~найпростішому зіткнювальному наближенні рівняння їхнього руху можна записати як
% ~~~~~~~~~~~~~~~~~~~~~~~~~~~~~~~~~~~~~~~~~~~~~~~~~~~~~~~~
\begin{equation}
	m_e\frac{d\vect v}{dt}
	=-e\vect E-\frac{m_e\vect v}{\tau_e},
\end{equation}
% ~~~~~~~~~~~~~~~~~~~~~~~~~~~~~~~~~~~~~~~~~~~~~~~~~~~~~~~~
де $\tau_e$~--- характерний час втрати впорядкованого руху електрона внаслідок зіткнень.

Для сталого поля
\begin{equation*}
	\vect v_d=-\frac{e\tau_e}{m_e}\vect E,
\end{equation*}
а~густина струму
% ~~~~~~~~~~~~~~~~~~~~~~~~~~~~~~~~~~~~~~~~~~~~~~~~~~~~~~~~
\begin{equation}
	\vect j=-n_e e\vect v_d
	=\sigma\vect E
\end{equation}
% ~~~~~~~~~~~~~~~~~~~~~~~~~~~~~~~~~~~~~~~~~~~~~~~~~~~~~~~~
має провідність
% ~~~~~~~~~~~~~~~~~~~~~~~~~~~~~~~~~~~~~~~~~~~~~~~~~~~~~~~~
\begin{equation}
	\tcbhighmath{\sigma=\frac{n_e e^2\tau_e}{m_e}}.
\end{equation}
% ~~~~~~~~~~~~~~~~~~~~~~~~~~~~~~~~~~~~~~~~~~~~~~~~~~~~~~~~

Оцінимо питому провідність повністю іонізованої плазми, що складається з електронів і позитивно заряджених іонів із зарядом $Ze$ кожний. Рух іонів, зважаючи на їхні більші маси, можна не враховувати й вважати, що весь ток створюється рухом легких електронів. Величина $\sigma$ визначається зіткненнями електронів з іонами. Зіткнення електронів між собою на величину тока не впливають, оскільки при таких зіткненнях повний імпульс електронів не змінюється. Від цих зіткнень можна відволіктися. Між іонами й електронами плазми діють \emphx{кулонівські сили притягання}. Це --- \emphx{далекодіючі сили}. Електрон порівняно рідко підходить до іона на такі малі відстані, щоб напрямок його руху змінився різко і мав характер стрибка. Набагато більше значення мають взаємодії електрона відразу з \emphx{дуже великою кількістю іонів}, при яких напрямок траєкторії електрона змінюється плавно й безперервно. Відхилення електрона на великі кути від першопочаткового напрямку руху відбуваються в результаті накопичення \emphx{малих відхилень} при взаємодії його з «далекими» іонами. Тому про зіткнення, довжину й час вільного пробігу можна говорити лише в умовному смислі. Умовимося називати часом вільного пробігу електрона проміжок часу $\tau$, протягом якого напрямок руху електрона змінюється на кут порядка $90^\circ$.

Для оцінки величини $\tau$ допустимо спочатку, що єдиний електрон рухається в полі позитивного іона із зарядом $Ze$. Якщо $v$ --- швидкість електрона на нескінченності, а $b$ --- прицільний параметр, то при проходженні мимо іона траєкторія електрона відхиляється на кут, який визначається формулою
% ~~~~~~~~~~~~~~~~~~~~~~~~~~~~~~~~~~~~~~~~~~~~~~~~~~~~~~~~
\begin{equation*}
	\ctg\frac{\vartheta}{2} = \frac{mv^2b}{Ze^2},
\end{equation*}
% ~~~~~~~~~~~~~~~~~~~~~~~~~~~~~~~~~~~~~~~~~~~~~~~~~~~~~~~~
де $m$ --- маса електрона. Прицільний параметр $b$, для якого $\vartheta = 90^\circ$, визначається виразом $b = Ze^2/(mv^2)$. Йому відповідає «ефективний поперечний переріз»
% ~~~~~~~~~~~~~~~~~~~~~~~~~~~~~~~~~~~~~~~~~~~~~~~~~~~~~~~~
\begin{equation*}
	\sigma = \pi b^2 = \pi\left(\frac{Ze^2}{mv^2}\right)^2.
\end{equation*}
% ~~~~~~~~~~~~~~~~~~~~~~~~~~~~~~~~~~~~~~~~~~~~~~~~~~~~~~~~

Врахування далеких взаємодій приводить до того ж результату, але збільшеному в $L$ разів:
% ~~~~~~~~~~~~~~~~~~~~~~~~~~~~~~~~~~~~~~~~~~~~~~~~~~~~~~~~
\begin{equation*}
	\sigma = \pi L\left(\frac{Ze^2}{mv^2}\right)^2.
\end{equation*}
% ~~~~~~~~~~~~~~~~~~~~~~~~~~~~~~~~~~~~~~~~~~~~~~~~~~~~~~~~

Коефіцієнт $L$ називається \emphx{кулонівським логарифмом}. Він майже не залежить від температури й густини плазми. Для плазми, що складається з повністю іонізованого дейтерію, при $kT \sim 10$~кеВ і концентрації електронів $n \sim 10^{12}$--$10^{15}$~см$^{-3}$ $L \approx 15$. Так як кожний позитивний іон містить $Z$ елементарних зарядів, то концентрація таких іонів буде $n/Z$, а середня довжина й час «вільного пробігу»
% ~~~~~~~~~~~~~~~~~~~~~~~~~~~~~~~~~~~~~~~~~~~~~~~~~~~~~~~~
\begin{equation*}
	\bar{l} = \frac{1}{\sigma n/Z} = \frac{Z}{\sigma n}, \quad \bar{\tau} \approx \frac{\bar{l}}{v} = \frac{m^2\bar{v}^3}{\pi Z n e^4 L}.
\end{equation*}
% ~~~~~~~~~~~~~~~~~~~~~~~~~~~~~~~~~~~~~~~~~~~~~~~~~~~~~~~~

Підставивши сюди $m\bar{v}^2 \approx 3kT$, отримаємо
% ~~~~~~~~~~~~~~~~~~~~~~~~~~~~~~~~~~~~~~~~~~~~~~~~~~~~~~~~
\begin{equation}\label{eq:relaxation-time}
	\bar{\tau} \approx \frac{(3kT)^{3/2}\sqrt{m}}{\pi Z n e^4 L}.
\end{equation}
% ~~~~~~~~~~~~~~~~~~~~~~~~~~~~~~~~~~~~~~~~~~~~~~~~~~~~~~~~

Звідси для питомої провідності плазми знаходимо
% ~~~~~~~~~~~~~~~~~~~~~~~~~~~~~~~~~~~~~~~~~~~~~~~~~~~~~~~~
\begin{equation}\label{eq:plasma-conductivity}
	\sigma \approx \frac{ne^2\bar{\tau}}{m} \approx \frac{(3kT)^{3/2}}{\pi Z e^2 L\sqrt{m}}.
\end{equation}
% ~~~~~~~~~~~~~~~~~~~~~~~~~~~~~~~~~~~~~~~~~~~~~~~~~~~~~~~~

Звичайно, наведений висновок треба розглядати як грубу оцінку, а не як доказ формули~\eqref{eq:plasma-conductivity}. Однак можна дати й більш обґрунтований висновок цієї формули. Підставляючи числові значення й виражаючи енергетичну температуру $\Theta = kT$ в електронвольтах, приходимо до оціночної формули для $\sigma$ (в обернених секундах):
% ~~~~~~~~~~~~~~~~~~~~~~~~~~~~~~~~~~~~~~~~~~~~~~~~~~~~~~~~
\begin{equation*}
	\sigma \approx \frac{10^{13}}{Z}\Theta^{3/2},
\end{equation*}
% ~~~~~~~~~~~~~~~~~~~~~~~~~~~~~~~~~~~~~~~~~~~~~~~~~~~~~~~~
причому ми прийняли $L \approx 15$. Провідність плазми росте пропорційно абсолютній температурі в степені $3/2$. У гарячій плазмі провідність стає дуже високою. Так, при $\Theta = 10$~кеВ для дейтерієвої плазми $\sigma \approx 10^{19}$~с$^{-1}$, тобто більше, ніж у міді ($5\cdot10^{17}$~с$^{-1}$). Ще швидше росте з температурою теплопровідність плазми, а саме пропорційно $\Theta^{5/2}$, так як для плазми, очевидно, повинен бути справедливий закон Відемана--Франца.

У~повністю іонізованій слабкозв'язаній плазмі кулонівські зіткнення мають далекодійний характер. Результат кінетичного розрахунку має характерну температурну залежність
% ~~~~~~~~~~~~~~~~~~~~~~~~~~~~~~~~~~~~~~~~~~~~~~~~~~~~~~~~
\begin{equation}
	\sigma\propto
	\frac{T_e^{3/2}}{\ln\Lambda},
\end{equation}
% ~~~~~~~~~~~~~~~~~~~~~~~~~~~~~~~~~~~~~~~~~~~~~~~~~~~~~~~~
де $\ln\Lambda$~--- кулонівський логарифм, що слабко залежить від параметрів плазми. Таким чином, на відміну від слабко іонізованого газу, провідність повністю іонізованої плазми швидко зростає з температурою.

Для базового курсу важливий не стільки числовий коефіцієнт цієї формули, скільки фізична причина результату: зі зростанням температури електрони рухаються швидше, але водночас ефективність кулонівського розсіяння зменшується.

Велика різниця в масах електронів і іонів плазми робить можливим у плазмі існування таких квазірівноважних станів, які у відомому наближенні можуть бути охарактеризовані \emphx{двома температурами}. Дійсно, нехай початковий розподіл швидкостей електронів і іонів плазми ізотропний, але не максвеллівський. При зіткненні електрона з іншим електроном вони обмінюються енергією, величина якої порядка початкової енергії самих електронів. Тому час встановлення розподілу електронів за енергіями (тобто максвеллівського розподілу) із-за зіткнень між ними можна оцінити за формулою~\eqref{eq:relaxation-time}. (Це стане ясно, якщо в ній масу електрона $m$ замінити приведеною масою $m/2$.) Це час, який називається \emphx{електронним часом релаксації} $\bar{\tau}_e$, пропорційно квадратному кореню з маси електрона: $\bar{\tau}_e \sim \sqrt{m_e}$. Точно так же визначається \emphx{іонний час релаксації}, за яке встигає встановлюватися розподіл за енергіями між іонами із-за зіткнень між ними: $\bar{\tau}_i \sim \sqrt{m_i}$. Не то буде при взаємодії електронів з іонами. Тут за формулою~\eqref{eq:relaxation-time} получається час порядка $\sqrt{m_e}$. Однак при кожному зіткненні швидка частинка передає повільній лише незначну долю своєї енергії. В середньому ця доля порядка $m_e/m_i$ від першопочаткової енергії швидкої частинки. Для вирівнювання енергій потрібно релаксаційний час $\bar{\tau}_{ie}$, в $m_i/m_e$ раз більший, ніж $\bar{\tau}_e$. Таким чином,
% ~~~~~~~~~~~~~~~~~~~~~~~~~~~~~~~~~~~~~~~~~~~~~~~~~~~~~~~~
\begin{equation}\label{eq:relaxation-times-ratio}
	\bar{\tau}_e : \bar{\tau}_i : \bar{\tau}_{ie} \sim 1 : \sqrt{\frac{m_i}{m_e}} : \frac{m_i}{m_e}.
\end{equation}
% ~~~~~~~~~~~~~~~~~~~~~~~~~~~~~~~~~~~~~~~~~~~~~~~~~~~~~~~~

Звідси випливає, що $\bar{\tau}_e \ll \bar{\tau}_i \ll \bar{\tau}_{ie}$. Якщо плазму надати самій собі, то спочатку встановиться максвеллівський розподіл швидкостей електронів, потім --- іонів. Виникне квазірівноважний стан, в якому електрони будуть мати температуру $T_e$, а іони --- температуру $T_i$. Взагалі кажучи, $T_e \neq T_i$. В цьому випадку плазму називають \emphx{неізотермічною} або \emphx{двотемпературною}. Потім у результаті обміну енергіями між електронами й іонами встановиться максвеллівський розподіл для всієї плазми, який характеризується загальною температурою електронів та іонів (\emphx{ізотермічна плазма}).

Коли плазма знаходиться в електричному полі, то в ній починає текти електричний ток і виділятися джоулеве тепло. При цьому енергію від поля отримують майже виключно електрони, як найбільш рухливі частинки. Іони нагріваються головним чином за рахунок енергії, яку вони отримують від «гарячих» електронів при кулонівських взаємодіях з ними. Так як останній процес відбувається порівняно повільно, то \emphx{температура електронів у плазмі виявляється вищою за температуру іонів}. Різниця між ними може бути вельми значною. Так, у позитивному стовпі тліючого розряду при тисках порядка 0,1 міліметрів ртутного стовпа температура електронів може досягати 50~000~°C і вище, тоді як температура іонів не перевищує декількох сотень градусів.

%% --------------------------------------------------------
\section{Рух зарядженої частинки в~магнітному полі}\label{sec:magnetic-field-motion}
%% --------------------------------------------------------

У~магнітному полі на заряджену частинку діє сила Лоренца
% ~~~~~~~~~~~~~~~~~~~~~~~~~~~~~~~~~~~~~~~~~~~~~~~~~~~~~~~~
\begin{equation}
	\vect F=\frac{q}{c}\,\vect v\times\vect B.
\end{equation}
% ~~~~~~~~~~~~~~~~~~~~~~~~~~~~~~~~~~~~~~~~~~~~~~~~~~~~~~~~
За однорідного поля складова швидкості, паралельна $\vect B$, не змінюється, а~перпендикулярна складова приводить до рівномірного руху по колу. Траєкторія частинки є гвинтовою лінією.

Циклотронна частота та ларморівський радіус мають вигляд
% ~~~~~~~~~~~~~~~~~~~~~~~~~~~~~~~~~~~~~~~~~~~~~~~~~~~~~~~~
\begin{equation}
	\tcbhighmath{\omega_c=\frac{|q|B}{mc}},
	\qquad
	\tcbhighmath{r_L=\frac{mc\,v_\perp}{|q|B}}.
\end{equation}
% ~~~~~~~~~~~~~~~~~~~~~~~~~~~~~~~~~~~~~~~~~~~~~~~~~~~~~~~~
Електрони мають значно менший ларморівський радіус, ніж іони тієї самої температури, через малу масу. Гвинтову траєкторію частинки та геометричний зміст $r_L$ показано на \ref{fig:larmor-motion}.

\begin{figure}[htbp]\centering
	\begin{tikzpicture}[>=latex,scale=0.85]
		\draw[->] (-4.2,0)--(4.4,0) node[right] {$z\parallel\vect B$};\draw[->] (0,-2.2)--(0,2.5) node[above] {$y$};
		\draw[thick,domain=-3.7:3.7,samples=180] plot (\x,{1.15*sin(220*\x/360)});\node at (2.3,1.55) {$r_L$};\node at (0,-1.75) {$\displaystyle\omega_c=\frac{|q|B}{mc}$};
	\end{tikzpicture}
	\caption{Гвинтовий рух зарядженої частинки в~однорідному магнітному полі. Уздовж поля частинка рухається рівномірно, а~поперек нього обертається з~циклотронною частотою.}\label{fig:larmor-motion}\end{figure}

У~сильному магнітному полі, коли
\begin{equation*}
	r_L\ll L
\end{equation*}
і~зіткнення не встигають істотно змінити орбіту, заряджена частинка сильно прив'язана до силової лінії магнітного поля. Це лежить в~основі магнітного утримання плазми.

%% --------------------------------------------------------
\subsection{Магнітний момент і~рух у~неоднорідному полі}
%% --------------------------------------------------------

Круговий рух зарядженої частинки можна розглядати як струмовий контур. Його магнітний момент за модулем
% ~~~~~~~~~~~~~~~~~~~~~~~~~~~~~~~~~~~~~~~~~~~~~~~~~~~~~~~~
\begin{equation}
	\mu=\frac{mv_\perp^2}{2B}.
\end{equation}
% ~~~~~~~~~~~~~~~~~~~~~~~~~~~~~~~~~~~~~~~~~~~~~~~~~~~~~~~~
За повільної зміни магнітного поля, коли поле за один ларморівський період змінюється мало, ця величина є наближеним адіабатичним інваріантом:
% ~~~~~~~~~~~~~~~~~~~~~~~~~~~~~~~~~~~~~~~~~~~~~~~~~~~~~~~~
\begin{equation}
	\tcbhighmath{\mu=\frac{mv_\perp^2}{2B}\approx\mathrm{const}.}
\end{equation}
% ~~~~~~~~~~~~~~~~~~~~~~~~~~~~~~~~~~~~~~~~~~~~~~~~~~~~~~~~

Повна кінетична енергія частинки в~стаціонарному магнітному полі зберігається:
% ~~~~~~~~~~~~~~~~~~~~~~~~~~~~~~~~~~~~~~~~~~~~~~~~~~~~~~~~
\begin{equation}
	\frac{mv_\parallel^2}{2}+\frac{mv_\perp^2}{2}
	=
	\frac{mv_\parallel^2}{2}+\mu B
	=\mathrm{const}.
\end{equation}
% ~~~~~~~~~~~~~~~~~~~~~~~~~~~~~~~~~~~~~~~~~~~~~~~~~~~~~~~~
Отже, в~області зростання $B$ зменшується $v_\parallel$. За певних умов вона стає нульовою, після чого частинка рухається назад. На цьому ґрунтується принцип магнітного дзеркала.

%% --------------------------------------------------------
\section{Магнітне утримання плазми}\label{sec:magnetic-confinement}\hypertarget{magnetic-confinement}{}
%% --------------------------------------------------------

%% --------------------------------------------------------
\subsection{Діамагнетизм плазми}
%% --------------------------------------------------------

При накладенні магнітного поля на плазму електрони починають здійснювати обертальний рух навколо напрямку силових ліній поля з~циклотронною частотою
% ~~~~~~~~~~~~~~~~~~~~~~~~~~~~~~~~~~~~~~~~~~~~~~~~~~~~~~~~
\begin{equation*}
	\omega_0 = -\frac{e}{m_e c}\vect{B}
\end{equation*}
% ~~~~~~~~~~~~~~~~~~~~~~~~~~~~~~~~~~~~~~~~~~~~~~~~~~~~~~~~
(див. рис.~\ref{pic:cyclotron-motion}).

Цей рух приводить до виникнення токів
% ~~~~~~~~~~~~~~~~~~~~~~~~~~~~~~~~~~~~~~~~~~~~~~~~~~~~~~~~
\begin{equation*}
	J = \frac{e\omega_0}{2\pi}.
\end{equation*}
% ~~~~~~~~~~~~~~~~~~~~~~~~~~~~~~~~~~~~~~~~~~~~~~~~~~~~~~~~

Якщо радіус орбіти рівен $r$, то виникає магнітний момент орбітального руху
% ~~~~~~~~~~~~~~~~~~~~~~~~~~~~~~~~~~~~~~~~~~~~~~~~~~~~~~~~
\begin{equation*}
	m = \frac{J}{c}S = \frac{e\omega_0}{2\pi c}\pi r^2.
\end{equation*}
% ~~~~~~~~~~~~~~~~~~~~~~~~~~~~~~~~~~~~~~~~~~~~~~~~~~~~~~~~

\begin{figure}[h!]\centering
	% \includegraphics[width=0.5\linewidth]{cyclotron-motion}
	\caption{При накладенні магнітного поля заряджені частинки починають обертатися навколо напрямку силових ліній магнітного поля, що приводить до появи магнітного моменту}
	\label{pic:cyclotron-motion}
\end{figure}

Нехай компонента швидкості частинки, перпендикулярна магнітному полю, є $v_\perp$. Тоді радіус орбіти $r = v_\perp/\omega_0$. Відповідно магнітний момент виявляється рівним
% ~~~~~~~~~~~~~~~~~~~~~~~~~~~~~~~~~~~~~~~~~~~~~~~~~~~~~~~~
\begin{equation*}
	m = \frac{e\omega_0}{2\pi c}\pi\left(\frac{v_\perp}{\omega_0}\right)^2 = \frac{e}{2c}\frac{v_\perp^2}{\omega_0}.
\end{equation*}
% ~~~~~~~~~~~~~~~~~~~~~~~~~~~~~~~~~~~~~~~~~~~~~~~~~~~~~~~~

Підставляючи сюди вираз для циклотронної частоти, знаходимо
% ~~~~~~~~~~~~~~~~~~~~~~~~~~~~~~~~~~~~~~~~~~~~~~~~~~~~~~~~
\begin{equation*}
	m = \frac{mv_\perp^2}{2B}.
\end{equation*}
% ~~~~~~~~~~~~~~~~~~~~~~~~~~~~~~~~~~~~~~~~~~~~~~~~~~~~~~~~

Перепишемо цей вираз у~векторному вигляді, маючи на увазі, що індуційований магнітний момент напрямлений проти магнітного поля:
% ~~~~~~~~~~~~~~~~~~~~~~~~~~~~~~~~~~~~~~~~~~~~~~~~~~~~~~~~
\begin{equation*}
	\vect{m} = -\frac{mv_\perp^2}{2B^2}\vect{B}.
\end{equation*}
% ~~~~~~~~~~~~~~~~~~~~~~~~~~~~~~~~~~~~~~~~~~~~~~~~~~~~~~~~

Якщо концентрація зарядів у~плазмі є $n$, то намагніченість виявиться рівною
% ~~~~~~~~~~~~~~~~~~~~~~~~~~~~~~~~~~~~~~~~~~~~~~~~~~~~~~~~
\begin{equation*}
	\vect{I} = n\vect{m} \uparrow\downarrow \vect{B}.
\end{equation*}
% ~~~~~~~~~~~~~~~~~~~~~~~~~~~~~~~~~~~~~~~~~~~~~~~~~~~~~~~~

Таким чином, магнітна сприйнятливість плазми, яка визначається співвідношенням $\vect{I} = \kappa\vect{H}$, від'ємна: $\kappa < 0$. Це означає, що плазма має діамагнітні властивості. Зокрема, вона повинна виштовхуватися з~області сильного поля. На цій властивості основана ідея магнітного утримання плазми: необхідно створити магнітні \enquote{подушки}, від яких плазма буде відштовхуватися, в~результаті чого вона буде утримуватися в~обмеженій області простору, не торкаючись стінок камери. Сформульована ідея запропонована А.~Д.~Сахаровим і~І.~Є.~Таммом у~1950~р.

%% --------------------------------------------------------
\subsection{Магнітне дзеркало та конус втрат}
%% --------------------------------------------------------

Розглянемо плазму в~неоднорідному полі, яке слабше в~центральній частині та сильніше на кінцях. Для частинки, що рухається від області $B_{\min}$ до області $B_{\max}$, із законів збереження енергії та магнітного моменту маємо
% ~~~~~~~~~~~~~~~~~~~~~~~~~~~~~~~~~~~~~~~~~~~~~~~~~~~~~~~~
\begin{equation}
	\frac{v_{\perp}^2}{B}=\mathrm{const},
	\qquad
	v_\parallel^2+v_\perp^2=\mathrm{const}.
\end{equation}
% ~~~~~~~~~~~~~~~~~~~~~~~~~~~~~~~~~~~~~~~~~~~~~~~~~~~~~~~~
Якщо $\theta$~--- кут між швидкістю частинки та напрямком поля в області $B_{\min}$, то умовою відбиття є
% ~~~~~~~~~~~~~~~~~~~~~~~~~~~~~~~~~~~~~~~~~~~~~~~~~~~~~~~~
\begin{equation}
	\tcbhighmath{
		\sin^2\theta
		\geq
		\frac{B_{\min}}{B_{\max}}
	}.
\end{equation}
% ~~~~~~~~~~~~~~~~~~~~~~~~~~~~~~~~~~~~~~~~~~~~~~~~~~~~~~~~
Частинки з~малим пітч-кутом не відбиваються і~виходять із пастки. Вони утворюють так званий \emph{конус втрат} у~просторі швидкостей. Геометрію магнітного дзеркала та напрям руху частинки схематично показано на \ref{fig:magnetic-mirror}.

Магнітне дзеркало добре ілюструє важливу ідею: магнітне поле не \enquote{тримає} заряджену частинку силою вздовж однорідної силової лінії. Утримання виникає через зміну магнітного моменту поперечного руху в~неоднорідному полі.

\begin{figure}[htbp]\centering
    \localinput{mirror.tikz}
	\caption{Принцип магнітного дзеркала: у~міру зростання $B$ поперечна енергія збільшується за рахунок поздовжньої, і~за певних умов частинка відбивається.}\label{fig:magnetic-mirror}
\end{figure}

%% --------------------------------------------------------
\subsection{Тороїдальні пастки}
%% --------------------------------------------------------

Інший спосіб уникнути втрат через торці~--- замкнути силові лінії магнітного поля. Так виникають тороїдальні магнітні системи. У~токамаках та стеллараторах плазма утримується всередині тороїдальної камери.

У~токамаках частина магнітного поля створюється струмом, що тече самою плазмою. У~стелараторах основну конфігурацію поля створюють зовнішні котушки. Конкретна стійкість плазми визначається не тільки силою поля, а~й~геометрією конфігурації, тиском плазми, струмами та колективними нестійкостями. Схему замкнених силових ліній у~такій системі наведено на \ref{fig:toroidal-trap}.

\begin{figure}[htbp]\centering
	\begin{tikzpicture}[>=latex,scale=0.9]
		\draw[thick] (0,0) ellipse (3.7 and 1.55);\draw[thick] (0,0) ellipse (2 and 0.75);\draw[thick,->] (0,1.55) arc[start angle=90,end angle=35,x radius=3.7,y radius=1.55];\draw[thick,->] (2,0.75) arc[start angle=90,end angle=25,x radius=2,y radius=0.75];
		\node at (0,0) {плазма};\node[align=center] at (0,-2.1) {замкнені магнітні силові лінії\\тороїдальна магнітна пастка};
	\end{tikzpicture}
	\caption{Схематичний принцип тороїдального магнітного утримання: силові лінії замкнені, тому немає відкритих торців магнітної системи.}\label{fig:toroidal-trap}\end{figure}

%% --------------------------------------------------------
\subsection{Тиск плазми та магнітне поле}
%% --------------------------------------------------------

У~найпростішому наближенні тиск слабкозв'язаної плазми є сумою тисків її компонентів:
% ~~~~~~~~~~~~~~~~~~~~~~~~~~~~~~~~~~~~~~~~~~~~~~~~~~~~~~~~
\begin{equation}
	p\simeq n_e k_{\mathrm B}T_e+n_i k_{\mathrm B}T_i.
\end{equation}
% ~~~~~~~~~~~~~~~~~~~~~~~~~~~~~~~~~~~~~~~~~~~~~~~~~~~~~~~~
Магнітне поле також має енергетичну густину
% ~~~~~~~~~~~~~~~~~~~~~~~~~~~~~~~~~~~~~~~~~~~~~~~~~~~~~~~~
\begin{equation}
	w_B=\frac{B^2}{8\pi}.
\end{equation}
% ~~~~~~~~~~~~~~~~~~~~~~~~~~~~~~~~~~~~~~~~~~~~~~~~~~~~~~~~
Тому природним безрозмірним параметром є відношення
% ~~~~~~~~~~~~~~~~~~~~~~~~~~~~~~~~~~~~~~~~~~~~~~~~~~~~~~~~
\begin{equation}
	\tcbhighmath{
		\beta=\frac{p}{B^2/(8\pi)}
		=\frac{8\pi p}{B^2}
	}.
\end{equation}
% ~~~~~~~~~~~~~~~~~~~~~~~~~~~~~~~~~~~~~~~~~~~~~~~~~~~~~~~~
За малої $\beta$ густина магнітної енергії значно перевищує тиск плазми; за $\beta\sim1$ ці величини порівнянні.

Параметр $\beta$ важливий для магнітного утримання, оскільки плазма чинить тиск, що прагне розширити її, тоді як магнітне поле може забезпечувати необхідне утримання та змінювати напрям руху заряджених частинок. Графічне порівняння тиску плазми з~густиною магнітної енергії наведено на \ref{fig:plasma-beta}.

\begin{figure}[htbp]\centering
	\begin{tikzpicture}[>=latex,scale=0.9]
		\draw[->] (-0.2,0)--(5.2,0) node[right] {$B$};\draw[->] (0,-0.2)--(0,3.8) node[above] {$w_B=B^2/(8\pi)$};\draw[thick,domain=0:4.6] plot (\x,{0.14*\x*\x});\draw[thick] (0,2.6)--(4.6,2.6);\node[above right] at (3.7,2.6) {$p$};\node[align=center] at (2.5,-1) {$\displaystyle\beta=\frac{8\pi p}{B^2}$};
	\end{tikzpicture}
	\caption{Порівняння газового тиску плазми $p$ з~густиною магнітної енергії $B^2/(8\pi)$. Параметр $\beta$ характеризує їхнє співвідношення.}\label{fig:plasma-beta}\end{figure}

%% --------------------------------------------------------
\subsection{Основне рівняння магнітної гідродинаміки плазми}
%% --------------------------------------------------------

Виділимо в~плазмі елемент об'єму $dV = dxdydz$ (рис.~\ref{pic:volume-element}).

\begin{figure}[h!]\centering
	% \includegraphics[width=0.6\linewidth]{volume-element}
	\caption{Елемент об'єму в~плазмі, на який діють сили газокінетичного тиску}
	\label{pic:volume-element}
\end{figure}

Позначимо газокінетичний тиск, що діє в~плазмі як $P(x, y, z)$. Тоді вздовж осі $x$ на елемент об'єму буде діяти сила
% ~~~~~~~~~~~~~~~~~~~~~~~~~~~~~~~~~~~~~~~~~~~~~~~~~~~~~~~~
\begin{equation*}
	df_x = [P(x, y, z) - P(x + dx, y, z)]dydz = -\frac{\partial P}{\partial x}dV.
\end{equation*}
% ~~~~~~~~~~~~~~~~~~~~~~~~~~~~~~~~~~~~~~~~~~~~~~~~~~~~~~~~

Аналогічно вздовж осей $y$ і~$z$ будуть діяти сили
% ~~~~~~~~~~~~~~~~~~~~~~~~~~~~~~~~~~~~~~~~~~~~~~~~~~~~~~~~
\begin{equation*}
	df_y = -\frac{\partial P}{\partial y}dV, \quad df_z = -\frac{\partial P}{\partial z}dV.
\end{equation*}
% ~~~~~~~~~~~~~~~~~~~~~~~~~~~~~~~~~~~~~~~~~~~~~~~~~~~~~~~~

У~векторному вигляді сила записується у~вигляді $d\vect{f} = -\grad P \cdot dV$.

Запишемо рівняння руху розглядуваного елемента об'єму. Нехай $\rho_m$~--- масова густина плазми, так що маса елемента об'єму рівна $dm = \rho_m dV$. Тоді згідно з~другим законом Ньютона маємо
% ~~~~~~~~~~~~~~~~~~~~~~~~~~~~~~~~~~~~~~~~~~~~~~~~~~~~~~~~
\begin{equation*}
	dm\frac{d\vect{v}}{dt} = -\grad P \cdot dV, \quad\text{або}\quad \rho_m\frac{d\vect{v}}{dt} = -\grad P.
\end{equation*}
% ~~~~~~~~~~~~~~~~~~~~~~~~~~~~~~~~~~~~~~~~~~~~~~~~~~~~~~~~

Отримане рівняння визначає поведінку плазми за відсутності яких-небудь інших сил, крім газокінетичного тиску.

Оскільки в~плазмі рухаються заряди, то вони створюють токи, які приводять до появи магнітних полей. Ці поля в~свою чергу створюють сили Ампера, діючі на токи.

Сила Ампера, діюча на елемент об'єму $dV$, в~якому присутній ток з~густиною $\vect{j}$, рівна
% ~~~~~~~~~~~~~~~~~~~~~~~~~~~~~~~~~~~~~~~~~~~~~~~~~~~~~~~~
\begin{equation*}
	d\vect{F} = \frac{1}{c}\vect{j}\times\vect{B}dV.
\end{equation*}
% ~~~~~~~~~~~~~~~~~~~~~~~~~~~~~~~~~~~~~~~~~~~~~~~~~~~~~~~~

З~урахуванням цього рівняння руху приймає вигляд
% ~~~~~~~~~~~~~~~~~~~~~~~~~~~~~~~~~~~~~~~~~~~~~~~~~~~~~~~~
\begin{equation*}
	\rho_m\frac{d\vect{v}}{dt} = -\grad P + \frac{1}{c}\vect{j}\times\vect{B}.
\end{equation*}
% ~~~~~~~~~~~~~~~~~~~~~~~~~~~~~~~~~~~~~~~~~~~~~~~~~~~~~~~~

Придамо отриманому рівнянню другий вид. Відповідно до рівняння Максвелла ток зв'язаний з~магнітним полем співвідношенням
% ~~~~~~~~~~~~~~~~~~~~~~~~~~~~~~~~~~~~~~~~~~~~~~~~~~~~~~~~
\begin{equation*}
	\vect{j} = \frac{c}{4\pi}\rot\vect{B}
\end{equation*}
% ~~~~~~~~~~~~~~~~~~~~~~~~~~~~~~~~~~~~~~~~~~~~~~~~~~~~~~~~
(ми знехтували токами зміщення, тобто вважаючи швидкості зміни полей малими). Підстановка цього виразу в~рівняння руху приводить до наступного рівнянню:
% ~~~~~~~~~~~~~~~~~~~~~~~~~~~~~~~~~~~~~~~~~~~~~~~~~~~~~~~~
\begin{equation}\label{eq:MHD-equation}
	\tcbhighmath{\rho_m\frac{d\vect{v}}{dt} = -\grad P + \frac{1}{4\pi}\rot\vect{B}\times\vect{B}}.
\end{equation}
% ~~~~~~~~~~~~~~~~~~~~~~~~~~~~~~~~~~~~~~~~~~~~~~~~~~~~~~~~

Дане рівняння дозволяє описати основні особливості поведінки плазми в~магнітному полі.

%% --------------------------------------------------------
\subsection{Пінч-ефект}
%% --------------------------------------------------------

Плазму можна створити з~допомогою електричного розряду. Далі виникає задача утримати утворений плазмовий шнур від розпливання, зумовленого існуючим у~ньому газокінетичним тиском. Цієї мети можна досягти, використовуючи різні прийоми, до числа яких відноситься \emphx{пінч-ефект}.

\emphx{Пінч-ефектом} (від англ. \emphx{pinch}~--- звуження, стискання) називається властивість електричного токового каналу зменшувати свій перетин під дією магнітного поля, породжуваного самим током.

Іншими словами, токовий плазмовий шнур може утримуватися від розпливання, і~навіть стискатися, завдяки власному магнітному полю.

Вперше пінч-ефект описав у~1934~р. У.~Беннет стосовно до потоків швидких заряджених частинок у~газорозрядній плазмі. Термін \enquote{пінч-ефект} запропонував Л.~Тонкс у~1937~р., який вивчав дуговий розряд.

Розрізняють два типи пінч-ефекту:
\begin{enumerate}
	\item $z$-пінч, коли ток тече вздовж плазмового шнура і~взаємодіє з~власним магнітним полем,
	\item $\theta$-пінч, коли на плазмовий шнур накладається поздовжнє магнітне поле, яке створює кільцеві токи, в~результаті взаємодії яких з~полем відбувається стискання шнура.
\end{enumerate}

На рис.~\ref{pic:pinch-types} проілюстровані ці випадки. Розглянемо докладніше $z$-пінч.

\begin{figure}[h!]\centering
	% \includegraphics[width=0.7\linewidth]{pinch-types}
	\caption{Зліва~--- $z$-пінч, справа~--- $\theta$-пінч}
	\label{pic:pinch-types}
\end{figure}

%% --------------------------------------------------------
\subsection{Рівновага $z$-пінча}
%% --------------------------------------------------------

Будемо вважати, що
\begin{enumerate}
	\item густина тока постійна по перетину плазмового шнура: $j(r) = \text{const}$, так що повний ток $J = \pi R^2 j$;
	\item концентрація частинок постійна: $n(r) = \text{const}$.
\end{enumerate}

Магнітне поле всередині шнура знаходиться з~допомогою теореми про циркуляцію магнітного поля:
% ~~~~~~~~~~~~~~~~~~~~~~~~~~~~~~~~~~~~~~~~~~~~~~~~~~~~~~~~
\begin{equation*}
	\oint\limits_{\Gamma} \vect{H} d\vect{l} = 2\pi r H = \frac{4\pi}{c}J(r).
\end{equation*}
% ~~~~~~~~~~~~~~~~~~~~~~~~~~~~~~~~~~~~~~~~~~~~~~~~~~~~~~~~

Тут в~якості контуру $\Gamma$ вибрана окружність радіуса $r$, а~$J(r) = j\cdot\pi r^2$~--- повний ток, що пронизує контур $\Gamma$ (рис.~\ref{pic:z-pinch-equilibrium}). Звідси знаходимо
% ~~~~~~~~~~~~~~~~~~~~~~~~~~~~~~~~~~~~~~~~~~~~~~~~~~~~~~~~
\begin{equation*}
	B = \frac{2\pi}{c}jr.
\end{equation*}
% ~~~~~~~~~~~~~~~~~~~~~~~~~~~~~~~~~~~~~~~~~~~~~~~~~~~~~~~~

Перепишемо цей вираз у~вигляді
% ~~~~~~~~~~~~~~~~~~~~~~~~~~~~~~~~~~~~~~~~~~~~~~~~~~~~~~~~
\begin{equation*}
	j = \frac{c}{2\pi r}B.
\end{equation*}
% ~~~~~~~~~~~~~~~~~~~~~~~~~~~~~~~~~~~~~~~~~~~~~~~~~~~~~~~~

Згідно з~основним рівнянням магнітної гідродинаміки умову рівноваги можна записати у~вигляді
% ~~~~~~~~~~~~~~~~~~~~~~~~~~~~~~~~~~~~~~~~~~~~~~~~~~~~~~~~
\begin{equation*}
	\grad P = \frac{1}{c}\vect{j}\times\vect{B}.
\end{equation*}
% ~~~~~~~~~~~~~~~~~~~~~~~~~~~~~~~~~~~~~~~~~~~~~~~~~~~~~~~~

\begin{figure}[h!]\centering
	% \includegraphics[width=0.6\linewidth]{z-pinch-equilibrium}
	\caption{До висновку умови рівноваги $z$-пінча. Контур $\Gamma$ показано штриховою лінією}
	\label{pic:z-pinch-equilibrium}
\end{figure}

Оскільки $\vect{j} \perp \vect{B}$, то, як видно з~рис.~\ref{pic:z-pinch-equilibrium}, вектор $\vect{j}\times\vect{B}$ напрямлений до осі шнура, причому
% ~~~~~~~~~~~~~~~~~~~~~~~~~~~~~~~~~~~~~~~~~~~~~~~~~~~~~~~~
\begin{equation*}
	(\vect{j}\times\vect{B})_r = -jB = -\frac{2\pi}{c}j^2r.
\end{equation*}
% ~~~~~~~~~~~~~~~~~~~~~~~~~~~~~~~~~~~~~~~~~~~~~~~~~~~~~~~~

Відповідно умова рівноваги приймає вигляд
% ~~~~~~~~~~~~~~~~~~~~~~~~~~~~~~~~~~~~~~~~~~~~~~~~~~~~~~~~
\begin{equation*}
	\frac{dP}{dr} = -\frac{2\pi}{c^2}j^2r \quad\Rightarrow\quad P + \frac{\pi}{c^2}j^2r^2 = \text{const}.
\end{equation*}
% ~~~~~~~~~~~~~~~~~~~~~~~~~~~~~~~~~~~~~~~~~~~~~~~~~~~~~~~~

Заміняючи в~цій рівності
% ~~~~~~~~~~~~~~~~~~~~~~~~~~~~~~~~~~~~~~~~~~~~~~~~~~~~~~~~
\begin{equation*}
	P = nkT, \quad j = \frac{c}{2\pi r}B,
\end{equation*}
% ~~~~~~~~~~~~~~~~~~~~~~~~~~~~~~~~~~~~~~~~~~~~~~~~~~~~~~~~
перепишемо умову рівноваги наступним чином:
% ~~~~~~~~~~~~~~~~~~~~~~~~~~~~~~~~~~~~~~~~~~~~~~~~~~~~~~~~
\begin{equation*}
	nkT(r) + \frac{B^2(r)}{4\pi} = \text{const}.
\end{equation*}
% ~~~~~~~~~~~~~~~~~~~~~~~~~~~~~~~~~~~~~~~~~~~~~~~~~~~~~~~~

Отриману рівність можна інтерпретувати таким чином, що сума газокінетичного ($nkT$) і~магнітного ($B^2/4\pi$) тисків повинна бути постійною в~об'ємі плазми.

Покладаючи на зовнішній границі плазмового шнура $T(R) = 0$, отримуємо звідси
% ~~~~~~~~~~~~~~~~~~~~~~~~~~~~~~~~~~~~~~~~~~~~~~~~~~~~~~~~
\begin{equation*}
	T(r) = \frac{1}{4\pi kn}\left[B^2(R) - B^2(r)\right].
\end{equation*}
% ~~~~~~~~~~~~~~~~~~~~~~~~~~~~~~~~~~~~~~~~~~~~~~~~~~~~~~~~

Цьому співвідношенню можна придати другий вид, скориставшись формулою
% ~~~~~~~~~~~~~~~~~~~~~~~~~~~~~~~~~~~~~~~~~~~~~~~~~~~~~~~~
\begin{equation*}
	B = \frac{2\pi}{c}jr = \frac{2J}{cR^2}r.
\end{equation*}
% ~~~~~~~~~~~~~~~~~~~~~~~~~~~~~~~~~~~~~~~~~~~~~~~~~~~~~~~~

Це дає
% ~~~~~~~~~~~~~~~~~~~~~~~~~~~~~~~~~~~~~~~~~~~~~~~~~~~~~~~~
\begin{equation*}
	T(r) = \frac{\pi}{knc^2}\left(\frac{J}{\pi R^2}\right)^2(R^2 - r^2).
\end{equation*}
% ~~~~~~~~~~~~~~~~~~~~~~~~~~~~~~~~~~~~~~~~~~~~~~~~~~~~~~~~

Зокрема, звідси випливає
% ~~~~~~~~~~~~~~~~~~~~~~~~~~~~~~~~~~~~~~~~~~~~~~~~~~~~~~~~
\begin{equation*}
	N_1 kT(0) = J^2/c^2,
\end{equation*}
% ~~~~~~~~~~~~~~~~~~~~~~~~~~~~~~~~~~~~~~~~~~~~~~~~~~~~~~~~
де введена погонна густина частинок плазмового шнура на його осі $N_1 = \pi R^2 n$.

Отримана рівність називається \emphx{умовою рівноваги Беннета} (частковий випадок).

%% --------------------------------------------------------
\subsection{Нестійкість $z$-пінча}
%% --------------------------------------------------------

Плазма~--- це об'єкт, в~якому спостерігаються різні нестійкості, які приводять до порушення її однорідності, форми. Нестійкості в~свою чергу обмежують час життя плазми. Дане обставина виявляється особливо критичним в~задачах термоядерного синтезу, коли плазму потрібно утримувати тривалий час, щоб активізувати реакцію.

До числа нестійкостей пінча відноситься виникнення і~розвиток перетяжок, як це проілюстровано на рис.~\ref{pic:pinch-instability}. Розглянемо дане явище на прикладі $z$-пінча.

\begin{figure}[h!]\centering
	% \includegraphics[width=0.6\linewidth]{pinch-instability}
	\caption{Нестійкість $z$-пінча~--- виникнення перетяжок плазмового шнура}
	\label{pic:pinch-instability}
\end{figure}

Механізм виникнення перетяжок можна пояснити наступним чином. Нехай в~якомусь місці шнура випадковим чином виникло звуження. Оскільки ток, що тече по шнуру, постійний, то густина тока зростає за законом $j = J/\pi R^2$. На поверхні шнура діє магнітне поле $B(R) = 2J/cR$. В~результаті зростає амперова сила, що стискає шнур:
% ~~~~~~~~~~~~~~~~~~~~~~~~~~~~~~~~~~~~~~~~~~~~~~~~~~~~~~~~
\begin{equation*}
	F_A = \frac{1}{c}jB(R) = \frac{1}{c}\frac{J}{\pi R^2}\frac{2J}{cR} = \frac{2J^2}{c^2\pi R^3}.
\end{equation*}
% ~~~~~~~~~~~~~~~~~~~~~~~~~~~~~~~~~~~~~~~~~~~~~~~~~~~~~~~~

Звідси випливає, що стискаюча дія магнітного поля не урівноважується газокінетичним тиском $P$, практично не змінним в~цьому місці. В~результаті відбувається подальше звуження перетяжки~--- розвивається нестійкість.

Для боротьби з~даною нестійкістю було запропоновано \enquote{вморозити} в~плазму поздовжнє магнітне поле $B_z^{(e)}$. Оскільки магнітний потік у~провідному середовищі зберігається:
% ~~~~~~~~~~~~~~~~~~~~~~~~~~~~~~~~~~~~~~~~~~~~~~~~~~~~~~~~
\begin{equation*}
	\Phi = \pi R^2 B_z^{(e)} = \text{const},
\end{equation*}
% ~~~~~~~~~~~~~~~~~~~~~~~~~~~~~~~~~~~~~~~~~~~~~~~~~~~~~~~~
то при звуженні шнура поле зростає:
% ~~~~~~~~~~~~~~~~~~~~~~~~~~~~~~~~~~~~~~~~~~~~~~~~~~~~~~~~
\begin{equation*}
	B_z^{(e)} \sim 1/R^2.
\end{equation*}
% ~~~~~~~~~~~~~~~~~~~~~~~~~~~~~~~~~~~~~~~~~~~~~~~~~~~~~~~~

Вклад цього поля в~густину енергії $\Delta U \sim (B_z^{(e)})^2 \sim 1/R^4$. Вважаємо систему замкнутою, так що робота поля здійснюється за його рахунок енергії:
% ~~~~~~~~~~~~~~~~~~~~~~~~~~~~~~~~~~~~~~~~~~~~~~~~~~~~~~~~
\begin{equation*}
	\delta A_{\text{поле}} = f\delta R = -\delta U \sim \delta R/R^5.
\end{equation*}
% ~~~~~~~~~~~~~~~~~~~~~~~~~~~~~~~~~~~~~~~~~~~~~~~~~~~~~~~~

Отже, $f \sim 1/R^5 > 0$, тобто виникаюча сила протидіє стискаючій пучок амперовій силі. Крім того, при зменшенні радіуса пучка ця сила за величиною росте швидше, ніж амперова сила ($F_A \sim 1/R^3$). Таким чином можна стабілізувати $z$-пінч.

%% --------------------------------------------------------
\section{Керований термоядерний синтез}\label{sec:fusion}
%% --------------------------------------------------------

Основний практичний інтерес, який представляє фізика плазми, пов'язаний з рішенням проблеми \emphx{керованого термоядерного синтезу}. Для того щоб у речовині почалися достатньо інтенсивні термоядерні реакції, його необхідно нагріти до температури в декілька кілоелектронвольт або десятків кілоелектронвольт, а при таких температурах всяка речовина знаходиться в стані плазми. Найбільш перспективними «робочими речовинами» для термоядерного реактора являються \emphx{ізотопи водороду}: \emphx{дейтерій} (D) і \emphx{тритій} (T). Термоядерну реакцію синтезу легше отримати не в чистому дейтерії, а в його суміші з тритієм. Повне кількість дейтерію в океанах $4\cdot10^{13}$ тонн, що еквівалентно енергії $10^{20}$~кВт$\cdot$рік (повна споживана на всьому земному шарі потужність становить $10^{10}$~кВт). Тритій, як сильно радіоактивний елемент, в природних умовах не зустрічається, а получається штучно. В будущих термоядерних реакторах расход тритія должен с ізбитком пополняться воспроизводством (регенерацією) його в результаті опромінення $^6$Li нейтронами, получающимися в самих термоядерних реакторах.

Так як термоядерні реакції повинні відбуватися порівняно плавно й повільно, то виникає необхідність достатньо тривалого удержання гарячої плазми в обмеженому об'ємі робочої камери й ізоляції її від стінок цієї камери. Для цього пропонується використовувати \emphx{магнітну термоізоляцію}, тобто поміщати плазму в сильне магнітне поле, яке перешкоджає іонам і електронам переміщатися в поперечному напрямку й уходити на стінки камери.

Необхідне вимога, якому повинен задовольняти всякий термоядерний реактор, полягає в тому, щоб енергія, яка виділяється в ядерних реакціях, з ізбитком компенсувала затрати енергії від зовнішніх джерел. Основними джерелами втрат енергії являються \emphx{гальмівне випромінювання} електронів при кулонівських зіткненнях останніх, а також \emphx{магнітогальмівне} (\emphx{циклотронне} або \emphx{бетатронне}) випромінювання, яке виникає внаслідок прискореного руху електронів у магнітному полі. Для самопідтримуючих термоядерних реакцій потрібно нагріти плазму до деякої «критичної» температури (50~кеВ для реактора на чистому дейтерії та 10~кеВ для реактора на рівнокомпонентній суміші дейтерію з тритієм). При цьому, як показує розрахунок, повинен виконуватися так званий критерій Лоусона: $n\tau > 10^{16}$~с/см$^3$ для реактора D--D і $n\tau > 10^{14}$~с/см$^3$ для реактора D--T. Тут $n$ --- концентрація іонів плазми (одного знака), а $\tau$ --- середній час удержання плазми.

Основна трудність, що стоїть на шляху створення керованого термоядерного синтезу, пов'язана з отриманням \emphx{спокійної}, або \emphx{стійкої}, \emphx{плазми}. Справа в тому, що із-за далекодіючого характеру кулонівських сил у плазмі відбуваються різні колективні процеси, наприклад мимоволі виникаючі \emphx{шуми й коливання}, які роблять плазму \emphx{нестійкою}. Основні зусилля при рішенні проблеми керованого термоядерного синтезу й направлені на подавлення цих нестійкостей.

Сучасний стан галузі характеризується значним прогресом у досягненні умов запалювання. На Національній установці запалювання (NIF) у 2022 році вперше було досягнуто наукового запалювання термоядерної реакції з позитивним виходом енергії. Міжнародний термоядерний експериментальний реактор (ITER) будується на півдні Франції і має на меті продемонструвати тривале горіння дейтерій-тритієвої плазми. Німецький стелларатор W7-X успішно демонструє переваги оптимізованої тороїдальної конфігурації для стабільного утримання плазми. Проєкт SPARC, що розробляється в MIT, використовує високотемпературні надпровідники для створення надзвичайно сильних магнітних полів, що дозволяє суттєво зменшити розміри реактора. Ці та інші установки відкривають шлях до практичного використання керованого термоядерного синтезу як джерела енергії.

%% --------------------------------------------------------
\section*{Підсумки}
%% --------------------------------------------------------

Плазма~--- це колективна система заряджених частинок, у~якій кулонівська взаємодія визначає поведінку речовини на масштабах, значно більших за міжчастинкові відстані. Її головні властивості можна зрозуміти через декілька характерних масштабів.

\begin{itemize}
	\item На макроскопічних масштабах плазма зазвичай є \emph{квазінейтральною}:
	      \begin{equation*}
		      \rho=\sum_s q_s n_s\simeq0.
	      \end{equation*}
	      Це не означає точного занулення густини заряду в~кожній точці.

	\item Дебаївська довжина
	      \begin{equation*}
		      \lambda_D^{-2}
		      =4\pi\sum_s\frac{n_s q_s^2}{k_{\mathrm B}T_s}
	      \end{equation*}
	      визначає характерний просторовий масштаб електростатичного екранування.

	\item Колективний опис має сенс, якщо в~дебаївській сфері міститься багато частинок:
	      \begin{equation*}
		      N_D=\frac{4\pi}{3}n\lambda_D^3\gg1.
	      \end{equation*}
	      Разом із вимогою малості дебаївської довжини порівняно з~розміром системи це дає типову ієрархію
	      \begin{equation*}
		      a\ll\lambda_D\ll L.
	      \end{equation*}

	\item Власна частота електронних плазмових коливань
	      \begin{equation*}
		      \omega_{pe}
		      =\sqrt{\frac{4\pi n_e e^2}{m_e}}
	      \end{equation*}
	      задає характерний часовий масштаб електронної відповіді.

	\item Магнітне поле змушує заряджені частинки рухатися по гвинтових траєкторіях із ларморівським радіусом
	      \begin{equation*}
		      r_L=\frac{mc\,v_\perp}{|q|B}.
	      \end{equation*}
	      Тому порівняння $r_L$ з~іншими просторовими масштабами є принципово важливим для магнітного опису плазми.

	\item Тиск плазми та магнітне поле характеризуються безрозмірним параметром
	      \begin{equation*}
		      \beta=\frac{8\pi p}{B^2}.
	      \end{equation*}
	      Разом із геометрією магнітного поля цей параметр визначає умови, за яких магнітне поле може суттєво впливати на макроскопічну динаміку плазми.

	\item Опис окремих частинок є лише початком. Наступними рівнями теорії є магнітна гідродинаміка та кінетична теорія плазми, які враховують просторову неоднорідність, хвилі, зіткнення та колективні нестійкості.
\end{itemize}

%% --------------------------------------------------------
\section*{Задачі}
%% --------------------------------------------------------

\begin{problem}\label{prb:PlasmaDebye}
	\emphx{Дебаївське екранування.} Плазма має електронну концентрацію $n_e=10^{12}~\mathrm{см}^{-3}$ і~температуру електронів $T_e=2~\mathrm{eV}$. Знайдіть електронну дебаївську довжину $\lambda_{De}$ та число електронів у~дебаївській сфері.
\end{problem}

\begin{problem}\label{prb:PlasmaCollective}
	\emphx{Умова колективності.} Для плазми з~концентрацією $n=10^{14}~\mathrm{см}^{-3}$ і температурою $T=10~\mathrm{eV}$ оцініть середню відстань між частинками $a\sim n^{-1/3}$ та дебаївську довжину. Порівняйте їх і~визначте, чи виконується умова $N_D\gg1$.
\end{problem}

\begin{problem}\label{prb:PlasmaFrequency}
	\emphx{Плазмова частота.} Знайдіть електронну плазмову частоту та період плазмових коливань для $n_e=10^{10}~\mathrm{см}^{-3}$. Якій звичайній частоті електромагнітної хвилі відповідає ця плазмова частота?
\end{problem}

\begin{problem}\label{prb:PlasmaIons}
	\emphx{Рух електронів та іонів.} Для квазінейтральної водневої плазми порівняйте внески електронів та протонів у~загальну плазмову частоту
	\begin{equation*}
		\omega_p^2=4\pi\sum_s\frac{n_s q_s^2}{m_s}.
	\end{equation*}
	Оцініть, на скільки відсотків відрізнятиметься повна частота від електронної, якщо $n_e=n_i$.
\end{problem}

\begin{problem}\label{prb:PlasmaWave}
	\emphx{Поширення електромагнітної хвилі.} Однорідна холодна плазма має концентрацію $n_e=10^{6}~\mathrm{см}^{-3}$. Знайдіть граничну частоту $\nu_p=\omega_{pe}/(2\pi)$, нижче якої електромагнітна хвиля в~простій моделі плазми не поширюється. Порівняйте її з частотою радіохвилі $1~\mathrm{MHz}$.
\end{problem}

\begin{problem}\label{prb:PlasmaConductivity}
	\emphx{Провідність плазми.} У~певному діапазоні температур для повністю іонізованої плазми вважайте $\ln\Lambda$ сталою, так що
	\begin{equation*}
		\sigma\propto T_e^{3/2}.
	\end{equation*}
	У~скільки разів зміниться провідність при нагріванні електронів від $10$ до $100~\mathrm{eV}$? Поясніть фізичний зміст цієї температурної залежності.
\end{problem}

\begin{problem}\label{prb:PlasmaLarmor}
	\emphx{Ларморівський радіус.} Електрон має кінетичну енергію $100~\mathrm{eV}$, а~магнітне поле дорівнює $10^4~\mathrm{G}$. Оцініть ларморівський радіус, вважаючи швидкість повністю перпендикулярною до поля. Порівняйте його з дебаївською довжиною плазми з~$n_e=10^{12}~\mathrm{см}^{-3}$ та $T_e=100~\mathrm{eV}$.
\end{problem}

\begin{problem}\label{prb:PlasmaMirror}
	\emphx{Магнітне дзеркало.} У~магнітному дзеркалі
	\begin{equation*}
		B_{\min}=2\times10^4~\mathrm{G},
		\qquad
		B_{\max}=8\times10^4~\mathrm{G}.
	\end{equation*}
	Знайдіть граничний пітч-кут частинки в~області слабкого поля. Для ізотропного розподілу швидкостей оцініть частку частинок, яка належить конусу втрат.
\end{problem}

\begin{problem}\label{prb:PlasmaBeta}
	\emphx{Тиск плазми та магнітне поле.} У~квазінейтральній дейтерієвій плазмі
	\begin{equation*}
		n_e=n_i=10^{14}~\mathrm{см}^{-3},
		\qquad
		T_e=T_i=10~\mathrm{keV}.
	\end{equation*}
	Знайдіть тиск плазми в~$\mathrm{дин/см^2}$. Яке магнітне поле в~гаусах відповідає значенню $\beta=1$?
\end{problem}

\begin{problem}\label{prb:PlasmaScales}
	\emphx{Порівняння масштабів.} Для плазми характерна ієрархія
	\begin{equation*}
		a\ll\lambda_D\ll L.
	\end{equation*}
	Поясніть фізичний зміст кожної з~двох нерівностей. Чому перша забезпечує статистичний характер екранування, а~друга дозволяє говорити про квазінейтральність на макроскопічних масштабах?
\end{problem}

\begin{problem}\label{prb:plasma-field}
	Напівпростір справа від нескінченної зарядженої площини $AB$ (рис.~\ref{pic:plasma-field}) заповнено повністю іонізованою плазмою, яка складається з однозарядних позитивних іонів і електронів. Знайти електричне поле й розподіл зарядів у плазмі на будь-яких відстанях від площини $AB$, якщо поверхнева густина зовнішніх зарядів на площині $AB$ рівна $\sigma$.
\end{problem}

\begin{problem}\label{prb:point-charge-plasma}
	Розв'язати попередню задачу для точкового заряду $q$, оточеного з усіх сторін плазмою. Обмежитися відстанями від заряду $r$, на яких виконується умова $|e\varphi| \ll kT$.
\end{problem}

\begin{problem}\label{prb:surface-layer-thickness}
	В електростатиці вважається, що в стані рівноваги електрику розподіляється по поверхні провідника. Фактично воно розподіляється не по математичній поверхні, а всередині поверхневого шару кінцевої товщини $l$. Оцінити цю товщину.
\end{problem}

%% --------------------------------------------------------
\section*{Контрольні питання}
%% --------------------------------------------------------

\begin{Questions}
	\begin{quizlist}
		\item Що називають плазмою? Чому простої наявності заряджених частинок недостатньо для визначення плазми як колективної системи?

		\item Що таке квазінейтральність? Чому вона не означає точного виконання $\rho=0$ у~кожній точці плазми?

		\item Чим відрізняються температури електронів та іонів? За яких умов можна очікувати $T_e=T_i$, а~коли можливе $T_e\ne T_i$?

		\item Що таке дебаївська довжина? Який фізичний процес визначає цей масштаб?

		\item Виведіть дебаївське екранування з~больцманівського розподілу електронів і~рівняння Пуассона. Які наближення використовуються під час виведення?

		\item Чим відрізняються електронна дебаївська довжина $\lambda_{De}$ та загальна дебаївська довжина $\lambda_D$? Коли внеском іонів можна знехтувати?

		\item Що таке дебаївська сфера? Який фізичний зміст має число
		\begin{equation*}
			N_D=\frac{4\pi}{3}n\lambda_D^3?
		\end{equation*}
		Чому для колективного опису потрібно $N_D\gg1$?

		\item Що означає нерівність
		\begin{equation*}
			a\ll\lambda_D\ll L?
		\end{equation*}
		Які фізичні масштаби порівнюються в~цій нерівності?

		\item Що таке плазмова частота? Поясніть фізичний механізм електронних плазмових коливань і~виведіть оцінку
		\begin{equation*}
			\omega_{pe}
			=\sqrt{\frac{4\pi n_e e^2}{m_e}}.
		\end{equation*}

		\item Чому електромагнітна хвиля з~$\omega<\omega_{pe}$ не поширюється в~простій холодній моделі однорідної плазми? Що означає cutoff у~цьому випадку?

		\item Як магнітне поле змінює рух зарядженої частинки? Що таке циклотронна частота та ларморівський радіус?

		\item За яких умов рух частинки можна вважати магнітно замагніченим? Чому порівняння $r_L$ з~характерним масштабом неоднорідності поля є важливим?

		\item Поясніть принцип магнітного дзеркала. Яка роль збереження магнітного моменту
		\begin{equation*}
			\mu=\frac{m v_\perp^2}{2B}?
		\end{equation*}

		\item Що таке конус втрат? Як змінюється його кут при збільшенні відношення $B_{\max}/B_{\min}$?

		\item У~чому полягає принцип тороїдального магнітного утримання? Чим принципово відрізняються токамак і стеларатор?

		\item Що характеризує параметр
		\begin{equation*}
			\beta=\frac{8\pi p}{B^2}?
		\end{equation*}
		Чому для аналізу магнітного утримання недостатньо знати лише величину магнітного поля?

		\item Які явища цього розділу вже неможливо адекватно описати моделлю незалежних частинок? Які наступні теоретичні рівні опису плазми ви знаєте?
	\end{quizlist}
\end{Questions}